% Options for packages loaded elsewhere
\PassOptionsToPackage{unicode}{hyperref}
\PassOptionsToPackage{hyphens}{url}
\PassOptionsToPackage{dvipsnames,svgnames,x11names}{xcolor}
\documentclass[
  11pt,
]{article}
\usepackage{xcolor}
\usepackage[margin=2.5cm]{geometry}
\usepackage{amsmath,amssymb}
\setcounter{secnumdepth}{5}
\usepackage{iftex}
\ifPDFTeX
  \usepackage[T1]{fontenc}
  \usepackage[utf8]{inputenc}
  \usepackage{textcomp} % provide euro and other symbols
\else % if luatex or xetex
  \usepackage{unicode-math} % this also loads fontspec
  \defaultfontfeatures{Scale=MatchLowercase}
  \defaultfontfeatures[\rmfamily]{Ligatures=TeX,Scale=1}
\fi
\usepackage{lmodern}
\ifPDFTeX\else
  % xetex/luatex font selection
\fi
% Use upquote if available, for straight quotes in verbatim environments
\IfFileExists{upquote.sty}{\usepackage{upquote}}{}
\IfFileExists{microtype.sty}{% use microtype if available
  \usepackage[]{microtype}
  \UseMicrotypeSet[protrusion]{basicmath} % disable protrusion for tt fonts
}{}
\makeatletter
\@ifundefined{KOMAClassName}{% if non-KOMA class
  \IfFileExists{parskip.sty}{%
    \usepackage{parskip}
  }{% else
    \setlength{\parindent}{0pt}
    \setlength{\parskip}{6pt plus 2pt minus 1pt}}
}{% if KOMA class
  \KOMAoptions{parskip=half}}
\makeatother
\usepackage{longtable,booktabs,array}
\newcounter{none} % for unnumbered tables
\usepackage{calc} % for calculating minipage widths
% Correct order of tables after \paragraph or \subparagraph
\usepackage{etoolbox}
\makeatletter
\patchcmd\longtable{\par}{\if@noskipsec\mbox{}\fi\par}{}{}
\makeatother
% Allow footnotes in longtable head/foot
\IfFileExists{footnotehyper.sty}{\usepackage{footnotehyper}}{\usepackage{footnote}}
\makesavenoteenv{longtable}
\usepackage{graphicx}
\makeatletter
\newsavebox\pandoc@box
\newcommand*\pandocbounded[1]{% scales image to fit in text height/width
  \sbox\pandoc@box{#1}%
  \Gscale@div\@tempa{\textheight}{\dimexpr\ht\pandoc@box+\dp\pandoc@box\relax}%
  \Gscale@div\@tempb{\linewidth}{\wd\pandoc@box}%
  \ifdim\@tempb\p@<\@tempa\p@\let\@tempa\@tempb\fi% select the smaller of both
  \ifdim\@tempa\p@<\p@\scalebox{\@tempa}{\usebox\pandoc@box}%
  \else\usebox{\pandoc@box}%
  \fi%
}
% Set default figure placement to htbp
\def\fps@figure{htbp}
\makeatother
% definitions for citeproc citations
\NewDocumentCommand\citeproctext{}{}
\NewDocumentCommand\citeproc{mm}{%
  \begingroup\def\citeproctext{#2}\cite{#1}\endgroup}
\makeatletter
 % allow citations to break across lines
 \let\@cite@ofmt\@firstofone
 % avoid brackets around text for \cite:
 \def\@biblabel#1{}
 \def\@cite#1#2{{#1\if@tempswa , #2\fi}}
\makeatother
\newlength{\cslhangindent}
\setlength{\cslhangindent}{1.5em}
\newlength{\csllabelwidth}
\setlength{\csllabelwidth}{3em}
\newenvironment{CSLReferences}[2] % #1 hanging-indent, #2 entry-spacing
 {\begin{list}{}{%
  \setlength{\itemindent}{0pt}
  \setlength{\leftmargin}{0pt}
  \setlength{\parsep}{0pt}
  % turn on hanging indent if param 1 is 1
  \ifodd #1
   \setlength{\leftmargin}{\cslhangindent}
   \setlength{\itemindent}{-1\cslhangindent}
  \fi
  % set entry spacing
  \setlength{\itemsep}{#2\baselineskip}}}
 {\end{list}}
\usepackage{calc}
\newcommand{\CSLBlock}[1]{\hfill\break\parbox[t]{\linewidth}{\strut\ignorespaces#1\strut}}
\newcommand{\CSLLeftMargin}[1]{\parbox[t]{\csllabelwidth}{\strut#1\strut}}
\newcommand{\CSLRightInline}[1]{\parbox[t]{\linewidth - \csllabelwidth}{\strut#1\strut}}
\newcommand{\CSLIndent}[1]{\hspace{\cslhangindent}#1}
\setlength{\emergencystretch}{3em} % prevent overfull lines
\providecommand{\tightlist}{%
  \setlength{\itemsep}{0pt}\setlength{\parskip}{0pt}}
\usepackage{bookmark}
\IfFileExists{xurl.sty}{\usepackage{xurl}}{} % add URL line breaks if available
\urlstyle{same}
% fallback for those not using the hyperref driver hyperxmp:
\makeatletter
\@ifundefined{xmpquote}{\newcommand{\xmpquote}[1]{#1}}{}
\makeatother
\hypersetup{
  pdftitle={SCORE: A Framework for Synthesizing Complete{,} Owned{,} Responsive Environments},
  pdfauthor={J.P. Kardolus, Babon Innovations B.V., Utrecht, The Netherlands},
  colorlinks=true,
  linkcolor={Maroon},
  filecolor={Maroon},
  citecolor={Blue},
  urlcolor={Blue},
  pdfcreator={LaTeX via pandoc}}

\title{SCORE: A Framework for Synthesizing Complete, Owned, Responsive
Environments}
\author{J.P. Kardolus, Babon Innovations B.V., Utrecht, The Netherlands}
\date{Working draft, 9 October 2026}

\begin{document}
\maketitle
\begin{abstract}
Games, films and robot training need 3D worlds that are complete,
editable and owned by the people who make them. World models and
one-shot 3D generators can produce a convincing place from a sentence or
a picture, but they give frames that no engine can open, or one fused
scene without separate objects or clear rights. Studios have always
built worlds in stages, from concept art and layout to modelling,
surfacing, sound and review. We present SCORE, a world generation
framework that keeps these stages and automates large parts of them with
open vision foundation models, procedural tools and coding agents, while
the creator chooses and reviews. Because every stage's output is kept,
the creator can rerun or fix any single stage, and because every model
allows commercial use, the creator holds full commercial rights. We
demonstrate SCORE on the game world of 2099 and plan three further
worlds of different kinds. Open models, working in checkable stages, can
build worlds that a creator owns and keeps editing. Code and paper are
available at https://github.com/Babon-Innovations-b-v/score (MIT
license).
\end{abstract}

\begin{center}\fbox{\begin{minipage}{0.94\linewidth}\setlength{\parskip}{4pt}\small

\textbf{Gap: teaser figure.} This figure will show four worlds made with
SCORE: the game world of 2099, a stylised underwater world, a simulation
world for robots, and a third-person fantasy world. Each will appear as
its chosen concept beside the finished world in its engine. It waits
until the four worlds are finished and accepted by their creators.

\end{minipage}}\end{center}

\section{Introduction}\label{introduction}

Interactive 3D worlds are the raw material of games, of film and of
embodied AI, where agents are trained and tested in simulators {[}1{]}.
Building such a world is still expensive. Every object has to be
modelled, given surfaces, placed, given collision and physical
properties, lit and given sound, all in one consistent style, and the
world has to stay editable while the project changes. Generative models
promise to take over much of this work, but the people who would use
them are wary of it: in GDC's 2026 survey, 36 percent of game workers
used generative AI at work, and 52 percent thought it was hurting the
industry {[}2{]}. Taking over the work is therefore not enough. A tool
for building worlds has to keep the creator the author of the result,
holding full commercial rights.

Two lines of research generate the world itself. The first, learned
world models, generates it as frames. Genie learns an interactive
environment from video {[}3{]}, and Genie 3 renders a world that a
person can move through in real time; its announcement lists a
constrained range of actions and a few minutes of continuous interaction
among its limits {[}4{]}. Nothing that these models produce can be
opened in a game engine or an editor, and their world exists only as far
as the player looks. The second line, 3D world generators, produces
geometry instead. WorldGen turns a text prompt into a traversable scene
for standard game engines {[}5{]}, and HY-World 2.0 turns text or a
picture into a single splat or mesh world {[}6{]}. Most such outputs,
however, are ``static monolithic assets with limited editability and
physical interaction'' {[}7{]}, and WorldAct and WorldSculpt address
this by recovering separate objects from them after the fact {[}7{]},
{[}8{]}. In both lines the world is generated in one shot. When part of
the result is wrong, the creator can regenerate the whole world or
repair it by hand, but cannot see which step caused the fault, and
cannot change that step alone.

A third line of work keeps the world explicit. Infinigen and Infinigen
Indoors build every asset from procedural rules, with materials as
generators of their own {[}9{]}, {[}10{]}. This gives real geometry and
full control, but only over the content that the rules describe.
Holodeck has a language model write layout constraints, which a solver
then satisfies with assets retrieved from a library {[}1{]}. SceneCraft,
recursive code world models, WorldClaw and AutoUE let coding agents
write scenes, or whole games, as programs {[}11{]}, {[}12{]}, {[}13{]},
{[}14{]}. LEGO-Anything shows what such agents need. When coding agents
rebuild a scene from a single picture, they start weakly, undo their own
progress while editing, and judge their own geometry unreliably; giving
them grounded tools, rather than training them further, closes much of
that gap {[}15{]}. All of these systems, however, either rebuild a scene
that already exists or draw their objects from a fixed library of
assets.

Meanwhile, each of the parts that a world needs can now be read out of
pictures by an open model. Vision foundation models learned the world
from images in the way that language models learned it from text
{[}16{]}. They can now turn a picture into a 3D shape aligned with its
pixels {[}17{]}, into a mask for each named object {[}18{]}, into depth
in metres {[}19{]}, {[}20{]}, into the parts of an object {[}21{]} or
into a human body {[}22{]}; related models turn a sentence into motion
{[}23{]} or into sound {[}24{]}. Most of these models pass on pictures,
masks, depth maps or meshes rather than their internal features, so a
coding agent can connect them to one another as it would connect tools.
What none of them provides is one consistent style across objects that
different models have made. The field names this as an open problem
{[}25{]}, {[}26{]}, and Hunyuan3D Studio addresses it only by restyling
the input picture {[}27{]}.

Studios have always built worlds in stages: concept art, layout, asset
lists, modelling, surfacing, sound and review. Each stage has a clear
output that a person can check before the next one starts, which makes
this way of working the natural basis for a world generation framework.
SCORE keeps these stages and automates large parts of them with the
models above, with procedural tools and with coding agents, while the
creator chooses and reviews. Its name states what it aims for. The
worlds it makes are \textbf{complete}: every object, surface, light and
sound exists wherever the player looks. They are \textbf{owned}: the
creator decides what the world is, edits it in the tools they already
use, and stays its author, holding full commercial rights, because every
model in the framework allows commercial use of its output. They are
\textbf{responsive}: the world answers to the physics engine that loads
it, because every object has real collision, mass and material, so it
behaves physically in any engine or simulator, under physics that each
world may set for itself, such as the low gravity of the Moon. And they
can be any kind of \textbf{environment}, real or made up. SCORE does not
compete with procedural tools, part splitters or world generators; it
takes them in as building blocks wherever their licences allow. It also
does not generate a world in one shot: the creator chooses the look, the
concept and the style of each place, and coding agents and models do the
work of the stages between those choices. To our knowledge, no other
openly available world generation work joins vision foundation models
into engine-ready scenes that are built from reference pictures, steered
by a person's choices at each step, and made only with models that allow
commercial use of their output.

This paper makes three contributions:

\begin{itemize}
\tightlist
\item
  \textbf{The SCORE framework:} staged, retraceable generation from a
  creator's choices to one canonical scene, which consists of a
  generated base layer and the creator's own edit layers (Section 2).
\item
  \textbf{The mechanisms that keep its worlds coherent and make faults
  cheap to fix:} deterministic checks before every stage that costs
  money, a parts check that decides how each object is built so that
  every real-world object becomes one model, a shared library of
  rule-based surfaces that keeps shape and surface apart, and review
  tools for the creator (Section 2).
\item
  \textbf{Evidence from four worlds of different kinds:} the time, the
  cost and the faults caught before spending at every stage, and an
  evaluation on a shared benchmark (Sections 3 and 4).
\end{itemize}

\section{Methods}\label{methods}

\begin{figure}
\centering
\includegraphics[width=1\linewidth,height=\textheight,keepaspectratio,alt={The stages on one place, the central hub of the game 2099. (a) The concept, drawn by a picture model from the creator's reference pictures and chosen by the creator. (b) The dimensioned plan in metres, with walkways and standing spots. (c) An excerpt of the scene inventory, which has 74 rows. (d) Close-up pictures of two objects. (e) One surface from the shared library of rule-based surfaces at three wear settings. (f) The assembled room in the game engine.}]{figures/fig-overview.jpg}
\caption{The stages on one place, the central hub of the game 2099. (a)
The concept, drawn by a picture model from the creator's reference
pictures and chosen by the creator. (b) The dimensioned plan in metres,
with walkways and standing spots. (c) An excerpt of the scene inventory,
which has 74 rows. (d) Close-up pictures of two objects. (e) One surface
from the shared library of rule-based surfaces at three wear settings.
(f) The assembled room in the game engine.}
\end{figure}

Figure 1 shows the stages on one place. The creator chooses a concept
(a), and a coding agent turns it into a dimensioned plan (b) and a scene
inventory (c). An open picture model draws a close-up of every object in
the inventory, and a closed one redraws those that fail a shape check
(d). Each object is then built in code or generated in 3D, its parts are
painted from the shared library of rule-based surfaces (e), and the
place is assembled in an engine (f). The output of every stage is kept,
so the creator can trace a weak result to the stage that caused it,
rerun or fix that stage, and regenerate the stages after it while
keeping their own edits. The sections below describe each stage with an
example.

\subsection{Inventory as a partition of the
concept}\label{inventory-as-a-partition-of-the-concept}

\begin{figure}
\centering
\includegraphics[width=1\linewidth,height=\textheight,keepaspectratio,alt={The inventory of the lab on the Moon base, built from its concept. Left: the concept, cut into twelve tiles. Right: examples of how each visible element maps to an inventory row, or is dropped with a written reason. Of the 47 elements in this concept, 40 were mapped to inventory rows and 7 were dropped with a written reason.}]{figures/fig-partition.jpg}
\caption{The inventory of the lab on the Moon base, built from its
concept. Left: the concept, cut into twelve tiles. Right: examples of
how each visible element maps to an inventory row, or is dropped with a
written reason. Of the 47 elements in this concept, 40 were mapped to
inventory rows and 7 were dropped with a written reason.}
\end{figure}

The inventory decides what exists in a place. A coding agent cuts the
concept into tiles and lists every visible element, and each element
either becomes a row (with its kind, anchor, size and count) or is
dropped with a written reason (Figure 2). Nothing is made or placed that
is not a row. After the build, the place is rendered from the concept's
own camera and compared with the concept; a place that is clearly
sparser than its concept fails this concept-density check.

\subsection{Close-ups: the open model
first}\label{close-ups-the-open-model-first}

Every object is built from a close-up: a picture of the object alone on
a plain background, drawn from its crop of the concept.
Qwen-Image-Edit-2511 {[}28{]}, an open model that runs on our own rented
cards, draws every close-up first. A shape check then decides, without a
person, whether the picture may go on to the 3D step. Its first half
measures the picture: a plain border all round, one object filling a
sensible share of the frame, and an outline in the proportions of the
inventory row's box. Its second half is an open vision-language model,
Qwen3.8-27B {[}29{]}, which sees the crop beside the close-up and
answers set questions (one object, a clean background, the whole object,
the same object, no parts missing or added, the right proportions); each
question is asked three times with different seeds and the majority
decides. Only a close-up that fails is redrawn by Nano Banana Pro
{[}30{]}, first from the crop alone and then with the whole concept
beside it, and its pictures pass the same check. The check leans towards
failing, because a wrong picture that passes costs a bad 3D model, while
a good one that fails costs only a picture from the closed model
(Appendix E).

\subsection{Code or generation: the parts
check}\label{code-or-generation-the-parts-check}

\begin{figure}
\centering
\includegraphics[width=1\linewidth,height=\textheight,keepaspectratio,alt={The parts check decides how each object is built. Top: the close-up picture of four objects. Bottom: the object as built. The notice board and the tall locker are built in code, because their code builds show every part of their close-ups. The microscope and the desk chair are generated with Pixal3D, because no code build shows all of their parts.}]{figures/fig-parts-check.jpg}
\caption{The parts check decides how each object is built. Top: the
close-up picture of four objects. Bottom: the object as built. The
notice board and the tall locker are built in code, because their code
builds show every part of their close-ups. The microscope and the desk
chair are generated with Pixal3D, because no code build shows all of
their parts.}
\end{figure}

Picture-to-3D models build chunky solids well, but flat panels, thin
beams and openings badly, while code builders are exact and straight but
show only the parts that the coding agent wrote. For each kind of
object, the parts check renders a code build from the front and compares
it with the close-up picture: if every part of the close-up is visible
in the build, the kind is built in code, and otherwise it is generated
from the close-up with Pixal3D {[}17{]} and made solid (Figure 3).
Composite objects are split in the inventory, so that every real-world
object becomes one model: a console, for example, becomes a desk
carrying separate monitors and keyboards.

\subsection{Shape and surface apart}\label{shape-and-surface-apart}

\begin{figure}
\centering
\includegraphics[width=1\linewidth,height=\textheight,keepaspectratio,alt={(a) Three surfaces from the shared library of rule-based surfaces (painted panel, bare steel and deck plate) at wear settings 0, 0.4 and 0.8. (b) The hub's tool board when each generated piece kept the colours of its own picture (left), and after every surface came from the library and each tool became its own model (right). Both shots use the same camera and the same build of the game.}]{figures/fig-surfaces.jpg}
\caption{(a) Three surfaces from the shared library of rule-based
surfaces (painted panel, bare steel and deck plate) at wear settings 0,
0.4 and 0.8. (b) The hub's tool board when each generated piece kept the
colours of its own picture (left), and after every surface came from the
library and each tool became its own model (right). Both shots use the
same camera and the same build of the game.}
\end{figure}

Models and code builders give shape only. Every part of every object
takes one surface from a shared library of rule-based surfaces, written
as ProcFunc functions {[}31{]} that are built on Infinigen's shaders
{[}9{]} and take their colours from the place's palette (Figure 4a). A
room has one wear setting, and wear is applied where it has a cause,
such as edges and the places where feet pass. On a generated model,
PartCrafter {[}21{]} splits the shape into parts, and an open
vision-language model, shown each part outlined on the close-up, names
the library surface it takes; the picture's colours are not kept. Figure
4b shows why: when each piece kept the colours of its own picture, every
piece brought its own rust and stains, and the room looked inconsistent.
Sound is attached in the same way, to surfaces, rooms and objects, and
is generated by MOSS-SoundEffect {[}24{]} (Appendix F).

\subsection{Terrain}\label{terrain}

\begin{figure}
\centering
\includegraphics[width=1\linewidth,height=\textheight,keepaspectratio,alt={The stages for large outdoor places, on the ground around the Moon base. (a) The dimensioned plan with the base's flat areas and paths. (b) A diorama picture drawn over the plan. (c) The heights that follow the diorama, shown shaded. (d) The skin painted over the plan. (e) The result in the engine, seen from a crater rim.}]{figures/fig-terrain.jpg}
\caption{The stages for large outdoor places, on the ground around the
Moon base. (a) The dimensioned plan with the base's flat areas and
paths. (b) A diorama picture drawn over the plan. (c) The heights that
follow the diorama, shown shaded. (d) The skin painted over the plan.
(e) The result in the engine, seen from a crater rim.}
\end{figure}

Large outdoor places follow the same pattern (Figure 5). A dimensioned
plan places the flat areas and paths the place needs (a), a diorama
picture drawn over the plan sets the landforms (b), and the heights
follow the picture (c). A skin is painted over the plan, with fine
relief from monocular depth {[}20{]} (d), rocks are single generated
models, and everything is built on the world's own curved ground (e).

\subsection{Checks before spending}\label{checks-before-spending}

\begin{figure}
\centering
\includegraphics[width=1\linewidth,height=\textheight,keepaspectratio,alt={A fault caught by the placement check. (a) In the first build of the hub, a conduit box covered the porthole. (b) In the current framework, the placement check found the box in front of the opening and moved it 0.5 m along the wall before anything was spent.}]{figures/fig-checks.jpg}
\caption{A fault caught by the placement check. (a) In the first build
of the hub, a conduit box covered the porthole. (b) In the current
framework, the placement check found the box in front of the opening and
moved it 0.5 m along the wall before anything was spent.}
\end{figure}

Deterministic checks run before every stage that costs money. They test
real geometry rather than relying on an agent's judgement, because a
fault costs nothing to fix in the plan but an evening of the creator's
time once the place is built. The checks look for light leaking out of
the room, pieces in front of openings or off their surface, doors that
do not separate their two sides, models that are open, too thin or
leaning, and surfaces off the palette (Appendix D). Figure 6 shows one
fault the placement check caught. A failed check sends the work back one
stage.

\subsection{The canonical scene}\label{the-canonical-scene}

Every place is exported as one OpenUSD stage {[}32{]} in which each
object carries its kind, its transform in metres, its collision, the
library surface of each part and that surface's sounds, with the
generated base layer under an edit layer that belongs to the creator, so
that an edit survives when the base is generated again (Figure 7).

\begin{figure}
\centering
\includegraphics[width=1\linewidth,height=\textheight,keepaspectratio,alt={The wreck, one outdoor place of the game 2099, as an OpenUSD stage. (a) The place in the game, from one of its fixed cameras. (b) The stage loaded in Blender and rendered from the same camera. (c) Blender's object outlines drawn over the game's picture. (d) The same camera after an edit in the edit layer, made here by a script, that moved the round cover and gave the capsule's painted hull another library surface, and after the base layer was generated again.}]{figures/fig-usd.jpg}
\caption{The wreck, one outdoor place of the game 2099, as an OpenUSD
stage. (a) The place in the game, from one of its fixed cameras. (b) The
stage loaded in Blender and rendered from the same camera. (c) Blender's
object outlines drawn over the game's picture. (d) The same camera after
an edit in the edit layer, made here by a script, that moved the round
cover and gave the capsule's painted hull another library surface, and
after the base layer was generated again.}
\end{figure}

A place's stage holds more than its inventory's objects. A scene record
adds what the place needs around them, each piece built by a code
builder and painted from the library: the room's shell and stairs, the
ground out to the horizon, water, the far backdrop, the neighbouring
places seen through open doorways (each referenced as its own stage),
every light with the sky and the exposure, and the cameras that the
place is judged from.

Before a place is accepted, its loose objects are dropped in a physics
simulation on the stage's own ground with their own triangles, and each
rest pose is written back into the layout. Settling only corrects: an
object that would turn or drift further than a small limit keeps its
laid pose and fails the resting check, which also fails anything
floating, tipping, sunk into the ground or lying inside another piece
(Appendix D). A large move is a fault of the layout, not of physics, and
the layout follows the creator's concept. On the wreck, 11 of 22 pieces
rested as first laid. The coordinating agent then laid the wreck out
again after its concept and cut shallow dents in the ground where pieces
had struck it, and all 22 rested; the creator has not yet reviewed this
new layout. On the old station, the check found a pile of sandbags lying
inside the lander's legs, and once the pile was moved, all 29 pieces
rested.

A place is complete only when its evidence says so. Coding agents
judging their own scenes from renders agree with a measured score on
geometry only 45.8 percent of the time, below chance {[}15{]}, and ours
called places finished that the creator then found unfinished. SCORE
therefore computes completeness instead of asking for it. A completion
gate passes a place only when every inventory row has its object on the
stage or a written reason for dropping it, nothing visible stands in for
a made piece, no real resting fault is left, every model the creator
picked is still the one shown, every check was run on the stage as it is
now, and the review page was rendered from that stage. Every check
answers pass, fail or unknown, and unknown blocks like fail, as in
LEGO-Anything's validator; a hook in the agent's harness refuses to end
a session that claims a place is done while its gate does not pass. The
resting check became a triage. Each object is judged by what holds it
up, read from its inventory row (the floor, a wall, the ceiling, the
object it stands on, a hook it hangs from, or the kit's own frame), with
SceneEval's guards against false alarms {[}33{]}: an overlap counts only
if it survives moving the object 5 mm away, and a standing object needs
four contact points with its weight inside them. Cables, cloths and
rubble may sink into what they lie on by an amount kept as data, and a
loose object that moves more than 0.2 m or turns more than 8 degrees
when dropped does not rest as laid, after SAGE's stability rule
{[}34{]}. The real faults are grouped by cause, each group with one
close-up aimed at its worst case. Over world 1's 17 places, the plain
check flagged 4,104 of 12,032 objects; the triage kept 139 objects with
real faults in 102 groups and marked the other 3,965 as false alarms,
most of them kit pieces that stand where the kit lays them. The drop
test ran on every place with loose pieces on its floor except the
square, where it did not finish in three hours. The creator's pick of a
model is written onto each of its objects as the file's hash, and an
export that would show anything else for a picked row is refused unless
the creator's decision to replace it is written down. Finally, each
object is rendered alone from eight directions and compared with its
close-up by its outline and by DINOv2 similarity {[}16{]}, and the least
similar are listed first: in the lab, the four lowest were a mouse drawn
black, a ceiling lamp without its lens, a scanner without its screen and
a wall monitor with a blank screen. The collision queries behind these
checks follow the interface of ProcFunc's scene tools {[}31{]}, whose
code is not yet released, built on python-fcl and trimesh, and one
annotation step gives every check per-object masks, depth, normals and
occlusion boundaries from any camera, the same labels a robotics world
needs.

A place's people are part of its stage. Its cast records who is there,
how many, where, doing what and why each is there, and nobody is placed
who is not in it. Each person becomes a skinned character (UsdSkel) that
plays a clip from its own starting point, in a layer of its own between
the creator's edit layer and the base. A crowd is one point instancer of
a cheap body, with a prototype for each clip, phase and palette of
clothes. On the leader's walk, UsdSkel skins the converted body to
within about a micrometre of the points the glTF file's own skinning
gives. The prologue's square holds the leader on the podium, a front row
of 24 people mixed from six kit builds, and a crowd of 10,000.

Characters are made by the framework's character maker, which has one
entry point. A short spec gives a name and a picture or a few words, and
the maker returns a rigged, clothed, textured and animated character as
a glTF file and as UsdSkel, with a far body for crowds and a set of
clips. A person's whole chain runs on one rented card. The A-pose
picture is the creator's own, or FLUX.2 klein 4B {[}35{]} draws it from
the words. SAM3DBody-cpp {[}36{]}, a C++ and ONNX runner of SAM 3D Body
{[}22{]}, reads the body off the picture, and the body is fitted to
SOMA-X through MHR. Kimodo {[}23{]} writes the clips from sentences. The
clothes are GarmentCode's sewing patterns {[}37{]}, draped on the body
by Blender's cloth simulation {[}38{]}. Hi3DGen {[}39{]} makes the hair
from a clay picture of the bald head, and the hair is laid on the scalp
as a shell. FLUX.2 klein base 4B {[}40{]} with a reference-depth LoRA
{[}41{]} draws the face on the head's own depth view. The chain then
builds the character, converts it to UsdSkel and renders it for review.
Rebuilt from her own picture, Nev, a crew member of the Moon base, took
11.2 minutes of an H100, and a new person made from one sentence took 21
minutes (Appendix G). Animals use the same entry point with a route that
does not assume a human body: Pixal3D makes the mesh from a close-up,
UniRig {[}42{]} rigs a quadruped, a fish gets a spine rig made in code,
and the clips are made in code. Two parts are weak. Blender's drapes are
rougher than those of GarmentCode's own cloth solver, a fork of NVIDIA
Warp that we cannot use because its licence allows non-commercial
research only. The animals' clips are made in code, so the dog's paws
twist and its feet slide (Limitations).

A place's sound is also part of its stage, in its own layer. Every sound
is a UsdMedia spatial audio prim with its file, its gain and how it
plays. The place's room tones loop without a position. An object's own
hum or fan loops where the object stands. Each library surface points at
its footstep, impact and scrape sounds, which an engine plays when
something happens on that surface. The levels, buses and reaches are
those the game 2099 played, now kept as the framework's sound catalogue.
On the hub, the layer holds 2 room tones, 23 object sounds and 9 surface
sounds. The review's walk and the demo's videos play what their camera
hears along its path: each sound at its gain over its distance from the
eye, panned by where it stands across the view.

What moves in a place moves in its stage as well, in a layer of its own:
OpenUSD time samples on the prims that the base already holds, at the
game's own speeds and lengths. The airlock's doors sink 2.8 m into the
floor at 3 m a second, and its warning light burns while the pump
empties the chamber and while the tanks fill it again, five seconds
each. The garage's big door slides apart at 2.5 m a second under a
warning beacon, and the hangar's roof leaves slide 11.2 m at 3 m a
second. How long a door stands open is the player's choice in the game,
so the motion record sets those pauses itself and says so. In the launch
view, the rocket lifts off 3 seconds after its engines light and climbs
at the game's rocket thrust of 24 m/s² against Earth's gravity. The game
never draws this, because the player rides the launch inside the
capsule, so the climb is an addition of the record, timed by the game's
launch sequence. Each motion also lists the sounds that the game plays
with it and when, and the sound layer plays them from the moving prim.
The places' vehicles and robots do not move yet: no place of world 1
shows the game's rover, drone or robots, and the street's car stands
still in the game as well.

A place's particle effects are a stage layer as well, and each emitter
is written twice. One prim carries the effect's parameters in the game's
numbers: its emitter shape, rate, lifetime, speed, spread, gravity, size
and opacity over its life, the wind it follows and its seed, so that an
engine can make the effect again with its own particle system. Under it,
the particles are baked as points over a loop that repeats exactly, so
any USD reader shows them without knowing those parameters. The bake
repeats the particle rules of Godot, the game's engine, with the
framework's own seeded draws, so the same record always gives the same
particles. The creator's rule for effects is to start from real footage
and to use many fine particles, because big or sparse ones read as
random dots, and the records keep the game's counts. On the calm
afternoon that the camp's scene record holds, 30,000 grains of 3 cm blow
along the ground on 25 cells round the camp, and two dust devils of
6,300 particles each wander downwind at 1.4 m/s; a storm variant of the
same layer holds 180,630 particles. The flat holds the wisp and the
breath of the cigarette smoked on its balcony, and the wreck the 40,000
grains that a supply rocket's engine throws across the landing pad.
Seven of the game's ten kinds of emitter are placed. The dust thrown by
a boot, a rover's wheel and a digger's scoop is not, because no place of
world 1 holds the walker or the machine that throws it.

The framework stands apart from the game it was first built for. Every
file a record names is in the repository or in one release of it: the
first world's made models (199 MB packed) and its 121 sounds (36 MB),
each listed with its checksum, its source and its licence. A check fails
when any code or data names a file of the game's own tree.

Some surfaces need more than a flat library colour, so the framework
writes its own shaders as UsdPreviewSurface materials over pictures it
bakes. Any USD reader can draw these. One example is the Earth in the
Moon's sky. Its continents, deserts, ice and cloud come from the game's
own noise shader, ported to numpy and baked once into a
longitude-latitude picture. The place's own sun lights the ball, so its
night side falls dark. The sun's disc hangs along each place's sun as
wide as the game draws it. The game's lamp, screen and grow-light
materials are library surfaces.

\begin{center}\fbox{\begin{minipage}{0.94\linewidth}\setlength{\parskip}{4pt}\small

\textbf{Gap: the scene in a game engine.} The same stage loaded in a
game engine, with its collision and surfaces, and the mass and friction
of each object written into it. It waits on the game engine adapter and
on the framework recording mass and friction.

\end{minipage}}\end{center}

\subsection{Review tools}\label{review-tools}

Each place ends with the creator's review. The review tools that SCORE
supplies are pages that show each stage's output side by side,
before-and-after shots from the same cameras, walkthroughs inside the
engine, and the result of every check. An optional world step can turn
the concept into a whole room that can be walked through, so that walls
and objects the concept does not show can be seen; nothing it produces
is shipped.

A review page is a static folder built only from the files the stages
wrote, with the assembled scene rendered by Blender from the place's
OpenUSD stage, on rented cards, from the cameras of its scene record and
along a walk through it (Figure 8). Each view is laid beside the game's
own shot from the same spot, and a view that is much darker than the
game's fails a brightness check, which is how lamps whose light fell off
too fast were found. A place with people shows each of them close and
the crowd wide, and its walk is drawn at consecutive moments of the
stage's time, so that the people move.

\begin{figure}
\centering
\includegraphics[width=1\linewidth,height=\textheight,keepaspectratio,alt={The wreck's review page, from two runs of the route over the same place. (a) One model: its close-up, its labelled parts (one colour per library surface) and its baked model, each from the first run and the rerun, drawn from the same camera. (b) The assembled scene from one fixed camera, before and after the rerun. (c) The parts check for each take, with what it caught in the rerun: the split of seven takes did not register against their pictures, so each was painted whole.}]{figures/fig-review.jpg}
\caption{The wreck's review page, from two runs of the route over the
same place. (a) One model: its close-up, its labelled parts (one colour
per library surface) and its baked model, each from the first run and
the rerun, drawn from the same camera. (b) The assembled scene from one
fixed camera, before and after the rerun. (c) The parts check for each
take, with what it caught in the rerun: the split of seven takes did not
register against their pictures, so each was painted whole.}
\end{figure}

\section{Results across worlds}\label{results-across-worlds}

We build four worlds with the same stages: the game world of 2099 (a
Moon base and its surroundings, a camp on Mars, and a prologue on
Earth), a stylised underwater world, a simulation world for robots, and
a third-person fantasy world.

\begin{center}\fbox{\begin{minipage}{0.94\linewidth}\setlength{\parskip}{4pt}\small

\textbf{Gap: Table 1, results per world and stage.} For each world, this
table will give the number of places, the wall-clock time, the cost in
GPU time and in picture calls, the number of models built in code and
generated, and the faults caught before spending at each stage against
those found afterwards. All 17 places of world 1 are built and their
records are in Appendix B, but they become results only when the creator
has reviewed the world, which has not happened since its first room;
worlds 2 to 4 have not been started.

\end{minipage}}\end{center}

\begin{center}\fbox{\begin{minipage}{0.94\linewidth}\setlength{\parskip}{4pt}\small

\textbf{Gap: figure of one place per world.} This figure will show one
representative place from each world beside its concept, rendered from
the concept's own camera. It waits on the same worlds; every place will
be shown in Appendix A.

\end{minipage}}\end{center}

\section{Evaluation}\label{evaluation}

\begin{center}\fbox{\begin{minipage}{0.94\linewidth}\setlength{\parskip}{4pt}\small

\textbf{Gap: LEGO-Bench.} This section will report SCORE's scores on
LEGO-Bench {[}15{]}, a benchmark of scenes rebuilt from pictures. It
waits until the benchmark's code is released and until SCORE's scene
export and its Blender adapter are built.

\end{minipage}}\end{center}

\subsection{With and without the agent
tools}\label{with-and-without-the-agent-tools}

The tools that check an agent's work (the completion gate, the resting
triage with its drop test, the pick lock and the render-and-compare,
Section 2) were tested on three places of world 1: a room (the workshop,
814 objects), an outdoor place (the street, 1,198 objects) and a dense
room (the hangar, 1,657 objects). Each arm started from the same
exported stage and the same inventory, and the same agent (Claude Opus
5.5 in Claude Code, headless) got the same prompt: make the place as
complete as it can in about 50 minutes, through the place's layout data
only. One arm had the framework as it was before these tools; the other
had them. Both arms were measured afterwards with the same instruments
(Table 1). Each arm ran once, so the table shows what happened, not an
average.

\textbf{Table 1.} One session per arm and place. Faults are the resting
triage's real faults left on the final stage, measured against the arm's
own data; placeholders by the placeholder check; rows are inventory rows
with no object and no written reason. Agent cost is the session's model
cost in dollars; cloud is what its rented machines cost.

{\def\LTcaptype{none} % do not increment counter
\begin{longtable}[]{@{}
  >{\raggedright\arraybackslash}p{(\linewidth - 18\tabcolsep) * \real{0.1000}}
  >{\raggedright\arraybackslash}p{(\linewidth - 18\tabcolsep) * \real{0.1000}}
  >{\raggedright\arraybackslash}p{(\linewidth - 18\tabcolsep) * \real{0.1000}}
  >{\raggedright\arraybackslash}p{(\linewidth - 18\tabcolsep) * \real{0.1000}}
  >{\raggedright\arraybackslash}p{(\linewidth - 18\tabcolsep) * \real{0.1000}}
  >{\raggedright\arraybackslash}p{(\linewidth - 18\tabcolsep) * \real{0.1000}}
  >{\raggedright\arraybackslash}p{(\linewidth - 18\tabcolsep) * \real{0.1000}}
  >{\raggedright\arraybackslash}p{(\linewidth - 18\tabcolsep) * \real{0.1000}}
  >{\raggedright\arraybackslash}p{(\linewidth - 18\tabcolsep) * \real{0.1000}}
  >{\raggedright\arraybackslash}p{(\linewidth - 18\tabcolsep) * \real{0.1000}}@{}}
\toprule\noalign{}
\begin{minipage}[b]{\linewidth}\raggedright
Place, arm
\end{minipage} & \begin{minipage}[b]{\linewidth}\raggedright
Real faults (start → end)
\end{minipage} & \begin{minipage}[b]{\linewidth}\raggedright
Placeholders
\end{minipage} & \begin{minipage}[b]{\linewidth}\raggedright
Rows missing
\end{minipage} & \begin{minipage}[b]{\linewidth}\raggedright
Minutes
\end{minipage} & \begin{minipage}[b]{\linewidth}\raggedright
Agent cost
\end{minipage} & \begin{minipage}[b]{\linewidth}\raggedright
Output tokens
\end{minipage} & \begin{minipage}[b]{\linewidth}\raggedright
Cloud
\end{minipage} & \begin{minipage}[b]{\linewidth}\raggedright
Exports
\end{minipage} & \begin{minipage}[b]{\linewidth}\raggedright
Gate said unknown
\end{minipage} \\
\midrule\noalign{}
\endhead
\bottomrule\noalign{}
\endlastfoot
Workshop, without & 18 → 13 & 0 → 0 & 1 → 0 & 36 & \$2.26 & 30.4k &
€0.85 & 7 & not used \\
Workshop, with & 18 → 0 & 0 → 0 & 1 → 0 & 26 & \$1.42 & 14.3k & €0.85 &
2 & 2 of 4 calls \\
Street, without & 11 → 0 & 0 → 0 & 0 → 0 & 38 & \$2.40 & 30.0k & €0.32 &
9 & not used \\
Street, with & 12 → 0 & 0 → 0 & 0 → 0 & 39 & \$1.21 & 10.4k & €0.76 & 1
& 0 of 5 calls \\
Hangar, without & 14 → 4 & 5 → 0 & 0 → 0 & 47 & \$2.01 & 26.6k & €0.00 &
7 & not used \\
Hangar, with & 14 → 9 & 5 → 0 & 0 → 0 & 32 & \$2.34 & 29.1k & €0.64 & 5
& 3 of 6 calls \\
\end{longtable}
}

The tools made building cheaper more often than they made the scene
better. On the workshop and the street, the arm with the tools used
about half the agent's output and cost, and exported the stage once or
twice instead of seven to nine times, because the triage named the few
real faults among hundreds of alarms. On the workshop it left no real
fault, but 13 of its 18 were cleared by writing what holds a piece up
into the inventory (the tools on the tool board hang from it; the arm
monitors are clamped to their bench), which the creator still has to
confirm; the other five were moved off the wall they stood in. On the
hangar the arm without the tools left fewer faults: the arm with them
spent its time replacing five placeholder slabs with the kit's own wall
panels and ribs and trying the drop test, which failed on a room with no
ground mesh (fixed after the run), and the nine faults it left are
models with stray parts below their feet that no pose can fix. In no arm
did the completion gate's hook have to refuse a claim, because no agent
claimed a place was complete; both arms reported what was left. Every
unknown answer of the gate was a check not yet run on the stage as it
then was, which is what unknown means; the agents ran the checks and
went on. Concept fit and the per-part surface check are not in the
table, because their tools are not built yet, and no row carries an
owner's pick, so the pick lock was never tested by an agent.

\begin{center}\fbox{\begin{minipage}{0.94\linewidth}\setlength{\parskip}{4pt}\small

\textbf{Gap: ablations.} This section will compare the framework with
and without the optional world step, and with the checks before spending
switched off, on the same places (the agent tools are compared above).
It will measure how much of the concept is covered, how dense the built
place is, the cost, the faults found after spending, and the creator's
review. It waits until every stage runs from the framework repository.

\end{minipage}}\end{center}

\begin{center}\fbox{\begin{minipage}{0.94\linewidth}\setlength{\parskip}{4pt}\small

\textbf{Gap: style coherence.} This section will measure whether a place
keeps one style, using rendered frames: the colour distance to the
palette for each kind of source, the spread of lightness, and the
density of ink lines, validated against the creator's reviews. The
measure is designed but not yet built.

\end{minipage}}\end{center}

\begin{center}\fbox{\begin{minipage}{0.94\linewidth}\setlength{\parskip}{4pt}\small

\textbf{Gap: robots in the simulation world.} This section will show a
rigged robot loaded together with the simulation world in a robotics
simulator, with masses, friction and joints written into the scene, and
with motion from the framework's motion models. It waits on world 3.

\end{minipage}}\end{center}

\section{Limitations}\label{limitations}

\begin{itemize}
\tightlist
\item
  \textbf{One creator, one world so far.} Every verdict on the results
  is one creator's, and only the first world has been built.
\item
  \textbf{Code builders do not scale well.} They give clean objects, but
  a coding agent has to write one for every kind of object.
\item
  \textbf{Places come out sparser than their concepts.} The
  concept-density check catches this but does not fix it.
\item
  \textbf{One stage still depends on a closed, paid picture model.} Nano
  Banana Pro, a closed and paid service, draws the concepts and every
  close-up that the open model's pictures fail. On the last 17 close-ups
  of world 1, none of the open model's pictures passed the shape check,
  so the closed model drew every close-up that was accepted. Its daily
  limits set the pace of world 1, and it can change or be retired, which
  limits reproducibility and leaves a public framework depending on a
  closed service for one stage. A cheaper model of the same family did
  not replace it (Appendix E).
\item
  \textbf{Parts do not follow finishes.} PartCrafter's parts follow
  shape, not paint, so a finish that does not have its own part, such as
  a gold foil patch on a white hull or a brown seat on a grey frame, is
  still lost when the part is painted with one surface.
\item
  \textbf{Objects are physical but not interactive.} Seats, terminals
  and other usable objects do not work in the output.
\item
  \textbf{Characters' clothes and animals' motion are rough.} Blender's
  drapes have more small crumples than Warp's; the work suit's trousers
  end about 5 cm higher and balloon, and the space suit gathers at the
  waist with sleeves that stand off the arm. The animals' clips are made
  in code: the dog's paws twist, its sit reads as a crouch and its feet
  slide by up to 5.4 cm.
\end{itemize}

\section{Future work}\label{future-work}

\begin{itemize}
\tightlist
\item
  Build the three further worlds: a stylised underwater world, a
  simulation world for robots, and a third-person fantasy world.
\item
  Add a game engine adapter for the OpenUSD scene, and write each
  object's mass and friction into it.
\item
  Make assets interactable, which robotics in particular needs.
\item
  Finish the animals and the drapes: fins that move, animal motion
  retargeted from a captured set whose licence has been checked, and
  Blender's drape brought closer to Warp's.
\item
  Make the open picture model's close-ups pass the shape check more
  often, so that the closed model is needed less.
\item
  Wrap open solvers as stages for worlds at other scales, starting with
  orbital dynamics through REBOUND {[}43{]}.
\item
  Evaluate SCORE on LEGO-Bench {[}15{]}.
\end{itemize}

\section{Conclusion}\label{conclusion}

SCORE builds complete, owned and responsive worlds from a creator's
choices, by joining open vision foundation models, procedural tools and
coding agents into fixed stages whose outputs are kept. On its first
world, the parts check, one model per real-world object and the shared
library of rule-based surfaces made a room consistent, and the
deterministic checks before every stage that costs money caught faults
before money was spent. Whether the framework holds across different
kinds of worlds is what the remaining worlds and the benchmark will
show.

\section*{References}\label{references}
\addcontentsline{toc}{section}{References}

\protect\phantomsection\label{refs}
\begin{CSLReferences}{0}{0}
\bibitem[\citeproctext]{ref-yang2024holodeck}
\CSLLeftMargin{{[}1{]} }%
\CSLRightInline{{Y. Yang \emph{et al.}}, {``Holodeck: Language guided
generation of {3D} embodied {AI} environments,''} in \emph{Proceedings
of the IEEE/CVF conference on computer vision and pattern recognition
(CVPR)}, 2024.}

\bibitem[\citeproctext]{ref-gdc2026survey}
\CSLLeftMargin{{[}2{]} }%
\CSLRightInline{D. Argüello, {``{36 percent of game workers use
generative {AI}, but half think it's bad for the industry}.''} Game
Developer, Feb. 2026.}

\bibitem[\citeproctext]{ref-bruce2024genie}
\CSLLeftMargin{{[}3{]} }%
\CSLRightInline{{J. Bruce, M. Dennis, A. Edwards, J. Parker-Holder,
\emph{et al.}}, {``{Genie}: Generative interactive environments,''}
\emph{arXiv preprint arXiv:2402.15391}, 2024.}

\bibitem[\citeproctext]{ref-genie3}
\CSLLeftMargin{{[}4{]} }%
\CSLRightInline{J. Parker-Holder and S. Fruchter, {``{Genie 3: A new
frontier for world models}.''} Google DeepMind blog, Aug. 2025.
Available:
\url{https://deepmind.google/blog/genie-3-a-new-frontier-for-world-models/}}

\bibitem[\citeproctext]{ref-wang2025worldgen}
\CSLLeftMargin{{[}5{]} }%
\CSLRightInline{{D. Wang, H. Jung, T. Monnier, K. Sohn, \emph{et al.}},
{``{WorldGen}: From text to traversable and interactive {3D} worlds,''}
\emph{arXiv preprint arXiv:2511.16825}, 2025.}

\bibitem[\citeproctext]{ref-hyworld2}
\CSLLeftMargin{{[}6{]} }%
\CSLRightInline{Tencent Hunyuan, {``{HY-World} 2.0.''} GitHub
repository, 2026. Available:
\url{https://github.com/Tencent-Hunyuan/HY-World-2.0}}

\bibitem[\citeproctext]{ref-hu2026worldact}
\CSLLeftMargin{{[}7{]} }%
\CSLRightInline{J. Hu, J. Guo, J. Cen, C. Yang, S. Li, and W. Shen,
{``{WorldAct}: Activating monolithic {3D} worlds into interactive-ready
object-centric scenes,''} \emph{arXiv preprint arXiv:2605.15843}, 2026.}

\bibitem[\citeproctext]{ref-niu2026worldsculpt}
\CSLLeftMargin{{[}8{]} }%
\CSLRightInline{{M. Niu, J. He, R. Yu, L. Fu, \emph{et al.}},
{``{WorldSculpt}: Generating compositional worlds from grounded
videos,''} \emph{arXiv preprint arXiv:2609.05416}, 2026.}

\bibitem[\citeproctext]{ref-raistrick2023infinigen}
\CSLLeftMargin{{[}9{]} }%
\CSLRightInline{A. Raistrick \emph{et al.}, {``Infinite photorealistic
worlds using procedural generation,''} \emph{arXiv preprint
arXiv:2306.09310}, 2023.}

\bibitem[\citeproctext]{ref-raistrick2024indoors}
\CSLLeftMargin{{[}10{]} }%
\CSLRightInline{A. Raistrick \emph{et al.}, {``{Infinigen Indoors}:
Photorealistic indoor scenes using procedural generation,''} in
\emph{Proceedings of the IEEE/CVF conference on computer vision and
pattern recognition (CVPR)}, 2024.}

\bibitem[\citeproctext]{ref-hu2024scenecraft}
\CSLLeftMargin{{[}11{]} }%
\CSLRightInline{{Z. Hu, A. Iscen, A. Jain, T. Kipf, \emph{et al.}},
{``{SceneCraft}: An {LLM} agent for synthesizing {3D} scene as {Blender}
code,''} in \emph{International conference on machine learning (ICML)},
2024.}

\bibitem[\citeproctext]{ref-li2026rcwm}
\CSLLeftMargin{{[}12{]} }%
\CSLRightInline{Z. Li, Y. Liao, and B. Zhu, {``Recursive code world
models: Building complex worlds through recursive scene programs,''}
\emph{arXiv preprint arXiv:2609.11499}, 2026.}

\bibitem[\citeproctext]{ref-guo2026worldclaw}
\CSLLeftMargin{{[}13{]} }%
\CSLRightInline{C. Guo, J. Li, Y. Li, and Z. Huang, {``{WorldClaw}:
Agentic {3D} open-world generation at scale,''} \emph{arXiv preprint
arXiv:2608.05248}, 2026.}

\bibitem[\citeproctext]{ref-yin2026autoue}
\CSLLeftMargin{{[}14{]} }%
\CSLRightInline{{L. Yin, W. Cheng, Z. Qin, T. Huang, \emph{et al.}},
{``{AutoUE}: Automated generation of {3D} games in {Unreal Engine} via
multi-agent systems,''} \emph{arXiv preprint arXiv:2603.07106}, 2026.}

\bibitem[\citeproctext]{ref-li2026lego}
\CSLLeftMargin{{[}15{]} }%
\CSLRightInline{X. Li \emph{et al.}, {``{LEGO-Anything}: Coding agents
for {3D} scene reconstruction,''} \emph{arXiv preprint
arXiv:2609.36380}, 2026.}

\bibitem[\citeproctext]{ref-oquab2023dinov2}
\CSLLeftMargin{{[}16{]} }%
\CSLRightInline{{M. Oquab, T. Darcet, T. Moutakanni, \emph{et al.}},
{``{DINOv2}: Learning robust visual features without supervision,''}
\emph{arXiv preprint arXiv:2304.07193}, 2023.}

\bibitem[\citeproctext]{ref-li2026pixal3d}
\CSLLeftMargin{{[}17{]} }%
\CSLRightInline{D.-Y. Li \emph{et al.}, {``{Pixal3D}: Pixel-aligned {3D}
generation from images,''} \emph{arXiv preprint arXiv:2605.10922},
2026.}

\bibitem[\citeproctext]{ref-carion2025sam3}
\CSLLeftMargin{{[}18{]} }%
\CSLRightInline{{N. Carion \emph{et al.}}, {``{SAM 3}: {Segment
Anything} with concepts,''} \emph{arXiv preprint arXiv:2511.16719},
2025.}

\bibitem[\citeproctext]{ref-wang2025moge2}
\CSLLeftMargin{{[}19{]} }%
\CSLRightInline{{R. Wang \emph{et al.}}, {``{MoGe-2}: Accurate monocular
geometry with metric scale and sharp details,''} \emph{arXiv preprint
arXiv:2507.02546}, 2025.}

\bibitem[\citeproctext]{ref-yang2024dav2}
\CSLLeftMargin{{[}20{]} }%
\CSLRightInline{{L. Yang, B. Kang, Z. Huang, \emph{et al.}}, {``{Depth
Anything V2},''} \emph{arXiv preprint arXiv:2406.09414}, 2024.}

\bibitem[\citeproctext]{ref-lin2025partcrafter}
\CSLLeftMargin{{[}21{]} }%
\CSLRightInline{Y. Lin \emph{et al.}, {``{PartCrafter}: Structured {3D}
mesh generation via compositional latent diffusion transformers,''} in
\emph{Advances in neural information processing systems (NeurIPS)},
2025.}

\bibitem[\citeproctext]{ref-yang2026sam3dbody}
\CSLLeftMargin{{[}22{]} }%
\CSLRightInline{{X. Yang, D. Kukreja, D. Pinkus, others, and K. Kitani},
{``{SAM 3D Body}: Robust full-body human mesh recovery,''} \emph{arXiv
preprint arXiv:2602.15989}, 2026.}

\bibitem[\citeproctext]{ref-rempe2026kimodo}
\CSLLeftMargin{{[}23{]} }%
\CSLRightInline{{D. Rempe, M. Petrovich, Y. Yuan, H. Zhang, X. B. Peng,
\emph{et al.}}, {``{Kimodo}: Scaling controllable human motion
generation,''} \emph{arXiv preprint arXiv:2603.15546}, 2026.}

\bibitem[\citeproctext]{ref-moss2026}
\CSLLeftMargin{{[}24{]} }%
\CSLRightInline{OpenMOSS-Team, {``{MOSS-SoundEffect} v2.0.''} Hugging
Face, 2026. Available:
\url{https://huggingface.co/OpenMOSS-Team/MOSS-SoundEffect-v2.0}}

\bibitem[\citeproctext]{ref-wu2026production}
\CSLLeftMargin{{[}25{]} }%
\CSLRightInline{J. Wu, Z. Lou, J. Liu, D. Du, C. Guo, and S. Guo,
{``From visual synthesis to interactive worlds: Toward production-ready
{3D} asset generation,''} \emph{arXiv preprint arXiv:2604.23629}, 2026.}

\bibitem[\citeproctext]{ref-yang2026flowscene}
\CSLLeftMargin{{[}26{]} }%
\CSLRightInline{Z. Yang \emph{et al.}, {``{FlowScene}: Style-consistent
indoor scene generation with multimodal graph rectified flow,''}
\emph{arXiv preprint arXiv:2603.19598}, 2026.}

\bibitem[\citeproctext]{ref-lei2025hunyuanstudio}
\CSLLeftMargin{{[}27{]} }%
\CSLRightInline{{B. Lei, Y. Li, X. Liu, \emph{et al.}}, {``{Hunyuan3D
Studio}: End-to-end {AI} pipeline for game-ready {3D} asset
generation,''} \emph{arXiv preprint arXiv:2509.12815}, 2025.}

\bibitem[\citeproctext]{ref-qwenimageedit}
\CSLLeftMargin{{[}28{]} }%
\CSLRightInline{Qwen Team, {``{Qwen-Image-Edit-2511} model card.''}
Hugging Face, 2025. Available:
\url{https://huggingface.co/Qwen/Qwen-Image-Edit-2511}}

\bibitem[\citeproctext]{ref-qwen38}
\CSLLeftMargin{{[}29{]} }%
\CSLRightInline{Qwen Team, {``{Qwen3.8-27B} model card.''} Hugging Face,
2026. Available: \url{https://huggingface.co/Qwen/Qwen3.8-27B}}

\bibitem[\citeproctext]{ref-nanobananapro}
\CSLLeftMargin{{[}30{]} }%
\CSLRightInline{Google, {``{Gemini 3 Pro Image ({`Nano Banana Pro'})
through the Gemini {API}}.''} 2026. Available:
\url{https://ai.google.dev/}}

\bibitem[\citeproctext]{ref-raistrick2026procfunc}
\CSLLeftMargin{{[}31{]} }%
\CSLRightInline{A. Raistrick \emph{et al.}, {``{ProcFunc}:
Function-oriented abstractions for procedural {3D} generation in
{Python},''} \emph{arXiv preprint arXiv:2604.26943}, 2026.}

\bibitem[\citeproctext]{ref-openusd}
\CSLLeftMargin{{[}32{]} }%
\CSLRightInline{Pixar Animation Studios and Alliance for OpenUSD,
{``{OpenUSD}: Universal scene description.''} 2026. Available:
\url{https://openusd.org/}}

\bibitem[\citeproctext]{ref-tam2025sceneeval}
\CSLLeftMargin{{[}33{]} }%
\CSLRightInline{H. I. I. Tam, H. I. D. Pun, A. T. Wang, A. X. Chang, and
M. Savva, {``{SceneEval}: Evaluating semantic coherence in
text-conditioned {3D} indoor scene synthesis,''} \emph{arXiv preprint
arXiv:2503.14756}, 2025.}

\bibitem[\citeproctext]{ref-xia2026sage}
\CSLLeftMargin{{[}34{]} }%
\CSLRightInline{H. Xia \emph{et al.}, {``{SAGE}: Scalable agentic {3D}
scene generation for embodied {AI},''} \emph{arXiv preprint
arXiv:2602.10116}, 2026.}

\bibitem[\citeproctext]{ref-flux2klein}
\CSLLeftMargin{{[}35{]} }%
\CSLRightInline{Black Forest Labs, {``{FLUX.2 {[}klein{]} 4B} model
card.''} Hugging Face, 2026. Available:
\url{https://huggingface.co/black-forest-labs/FLUX.2-klein-4B}}

\bibitem[\citeproctext]{ref-sam3dbodycpp}
\CSLLeftMargin{{[}36{]} }%
\CSLRightInline{A. Qammaz, {``{SAM3DBody-cpp}: A standalone {C++}
inference engine for {SAM 3D Body}.''} GitHub, 2026. Available:
\url{https://github.com/AmmarkoV/SAM3DBody-cpp}}

\bibitem[\citeproctext]{ref-korosteleva2023garmentcode}
\CSLLeftMargin{{[}37{]} }%
\CSLRightInline{M. Korosteleva and O. Sorkine-Hornung, {``{GarmentCode}:
Programming parametric sewing patterns,''} \emph{ACM Transactions on
Graphics}, vol. 42, no. 6, p. 197, 2023.}

\bibitem[\citeproctext]{ref-blender}
\CSLLeftMargin{{[}38{]} }%
\CSLRightInline{Blender Foundation, {``{Blender}.''} 2026. Available:
\url{https://www.blender.org/}}

\bibitem[\citeproctext]{ref-ye2025hi3dgen}
\CSLLeftMargin{{[}39{]} }%
\CSLRightInline{C. Ye \emph{et al.}, {``{Hi3DGen}: High-fidelity {3D}
geometry generation from images via normal bridging,''} in
\emph{Proceedings of the IEEE/CVF international conference on computer
vision (ICCV)}, 2025.}

\bibitem[\citeproctext]{ref-flux2kleinbase}
\CSLLeftMargin{{[}40{]} }%
\CSLRightInline{Black Forest Labs, {``{FLUX.2 {[}klein{]} Base 4B} model
card.''} Hugging Face, 2026. Available:
\url{https://huggingface.co/black-forest-labs/FLUX.2-klein-base-4B}}

\bibitem[\citeproctext]{ref-refcontrol}
\CSLLeftMargin{{[}41{]} }%
\CSLRightInline{{thedeoxen}, {``{refcontrol}: Reference-depth {LoRA} for
{FLUX.2 {[}klein{]} 4B}.''} Hugging Face, 2026. Available:
\url{https://huggingface.co/thedeoxen/refcontrol-FLUX.2-klein-4B-reference-depth-lora}}

\bibitem[\citeproctext]{ref-zhang2025unirig}
\CSLLeftMargin{{[}42{]} }%
\CSLRightInline{J.-P. Zhang, C.-F. Pu, M.-H. Guo, Y.-P. Cao, and S.-M.
Hu, {``One model to rig them all: Diverse skeleton rigging with
{UniRig},''} \emph{ACM Transactions on Graphics (SIGGRAPH)}, 2025.}

\bibitem[\citeproctext]{ref-rein2012rebound}
\CSLLeftMargin{{[}43{]} }%
\CSLRightInline{H. Rein and S.-F. Liu, {``{REBOUND}: An open-source
multi-purpose {N}-body code for collisional dynamics,''} \emph{Astronomy
\& Astrophysics}, vol. 537, p. A128, 2012.}

\bibitem[\citeproctext]{ref-simeoni2025dinov3}
\CSLLeftMargin{{[}44{]} }%
\CSLRightInline{{O. Siméoni, H. V. Vo, M. Seitzer, \emph{et al.}},
{``{DINOv3},''} \emph{arXiv preprint arXiv:2508.10104}, 2025.}

\bibitem[\citeproctext]{ref-li2026segvigen}
\CSLLeftMargin{{[}45{]} }%
\CSLRightInline{L. Li \emph{et al.}, {``{SegviGen}: Repurposing {3D}
generative model for part segmentation,''} \emph{arXiv preprint
arXiv:2603.16869}, 2026.}

\bibitem[\citeproctext]{ref-deng2025geosam2}
\CSLLeftMargin{{[}46{]} }%
\CSLRightInline{K. Deng \emph{et al.}, {``{GeoSAM2}: Unleashing the
power of {SAM2} for {3D} part segmentation,''} \emph{arXiv preprint
arXiv:2508.14036}, 2025.}

\bibitem[\citeproctext]{ref-marble_terms}
\CSLLeftMargin{{[}47{]} }%
\CSLRightInline{World Labs, {``{Terms of Service}.''} 2026. Available:
\url{https://docs.worldlabs.ai/terms-of-service}}

\bibitem[\citeproctext]{ref-behnamghader2024llm2vec}
\CSLLeftMargin{{[}48{]} }%
\CSLRightInline{P. BehnamGhader, V. Adlakha, M. Mosbach, D. Bahdanau, N.
Chapados, and S. Reddy, {``{LLM2Vec}: Large language models are secretly
powerful text encoders,''} in \emph{Conference on language modeling
(COLM)}, 2024.}

\bibitem[\citeproctext]{ref-grattafiori2024llama3}
\CSLLeftMargin{{[}49{]} }%
\CSLRightInline{{A. Grattafiori, A. Dubey, A. Jauhri, \emph{et al.}},
{``The {Llama 3} herd of models,''} \emph{arXiv preprint
arXiv:2407.21783}, 2024.}

\bibitem[\citeproctext]{ref-claudecode}
\CSLLeftMargin{{[}50{]} }%
\CSLRightInline{Anthropic, {``{Claude Code}.''} 2026. Available:
\url{https://www.anthropic.com/claude-code}}

\end{CSLReferences}

\section*{Appendix}\label{appendix}
\addcontentsline{toc}{section}{Appendix}

\appendix

\section{Per-world demos}\label{per-world-demos}

\subsection{World 1: the game world of
2099}\label{world-1-the-game-world-of-2099}

Figure A1 shows one frame of each of world 1's 17 places, in story
order: the prologue on Earth, the Moon base and the ground around it,
and the first camp on Mars. Every frame is rendered by Blender's Cycles
from the place's OpenUSD stage, from one of the player's spots in its
scene record, with the stage's own lights and sky and its characters; no
game engine is involved. What a scene record lists as drawn only by the
game, such as the ink outline pass, the haze and the stars, is missing.
A walk-through video of every place was rendered the same way, two to
four slow shots of 3.5 seconds a place from the player's spots, and
joined in story order into one video of world 1, 3 minutes 35 seconds
long, with each place's plain name on a title card. For the videos,
every surface map was reduced to at most 1,024 pixels. The creator has
not yet reviewed world 1, so these pictures show the framework's output
as it stands, not an accepted result.

\begin{figure}
\centering
\includegraphics[width=1\linewidth,height=\textheight,keepaspectratio,alt={World 1 of SCORE, one frame of each of its 17 places in story order: the prologue on Earth (flat to launch view), the Moon base, its walkway tube and the ground around it (hub to old station), and the first camp on Mars. Rendered by Blender's Cycles from each place's OpenUSD stage, from one of the player's spots, with the stage's own lights, sky and characters; not reviewed by the creator. The habitat's and the flat's style texts came from the optional world step (Generated using World Labs).}]{figures/fig-world1.jpg}
\caption{World 1 of SCORE, one frame of each of its 17 places in story
order: the prologue on Earth (flat to launch view), the Moon base, its
walkway tube and the ground around it (hub to old station), and the
first camp on Mars. Rendered by Blender's Cycles from each place's
OpenUSD stage, from one of the player's spots, with the stage's own
lights, sky and characters; not reviewed by the creator. The habitat's
and the flat's style texts came from the optional world step (Generated
using World Labs).}
\end{figure}

The 17 places' videos were rendered as 5,292 frames at 1280 by 720
pixels, in 34 jobs spread over 20 rented L4 cards, in 78 minutes for
€15.82.

Figure A2 follows two objects through the stages, from the concept to
the place.

\begin{figure}
\centering
\includegraphics[width=1\linewidth,height=\textheight,keepaspectratio,alt={Two objects through the stages: the lab's desk chair (top) and the wreck's forward section (bottom). From left: the concept the creator chose, the close-up drawn from it, the generated model with its parts, the model painted from the surface library, and the model in its place on the stage. The forward section shows a known fault: its scorched paint has no part of its own and is lost in painting (Section 5).}]{figures/fig-stages.jpg}
\caption{Two objects through the stages: the lab's desk chair (top) and
the wreck's forward section (bottom). From left: the concept the creator
chose, the close-up drawn from it, the generated model with its parts,
the model painted from the surface library, and the model in its place
on the stage. The forward section shows a known fault: its scorched
paint has no part of its own and is lost in painting (Section 5).}
\end{figure}

\begin{center}\fbox{\begin{minipage}{0.94\linewidth}\setlength{\parskip}{4pt}\small

\textbf{Gap: worlds 2 to 4, and a public link to the videos.} This
appendix will show every place of the other three worlds in the same
way. The videos of world 1 are not yet published at a public address.

\end{minipage}}\end{center}

\section{Detailed per-stage results}\label{detailed-per-stage-results}

Table B1 gives the records of world 1's 17 places as the build logs and
the cloud ledger hold them. They are records, not yet results: the
creator has not reviewed world 1 since its first room, and worlds 2 to 4
have not been started. The cloud costs are lower bounds. The ledger
attributes a batch to a place by the work folders and take names in it;
review renders (€24.24), the shape-check judge (€22.81), the game's
shared test machines (€32.63), settling (€4.43), the open picture models
(€10.27) and the repainting and spreading work (€8.44) could not be
attributed to one place. No time or cost is recorded for each stage of a
place separately, only for each batch, so the table gives none. The
resting check is the one on each place's review page of 8 October,
before the night's complete stages were exported; it counts every piece
it fails, including many deck plates that it finds sunk by 4 cm, which
have not been judged as real faults or not.

\textbf{Table B1.} World 1, per place. Models are those built in code
and those generated. Cloud is the ledger's rented machine time
attributed to the place, in euros; pictures are the picture model's
calls in dollars, from the build logs. Faults are those the checks
caught before money was spent and those found after.

{\def\LTcaptype{none} % do not increment counter
\begin{longtable}[]{@{}
  >{\raggedright\arraybackslash}p{(\linewidth - 14\tabcolsep) * \real{0.1250}}
  >{\raggedright\arraybackslash}p{(\linewidth - 14\tabcolsep) * \real{0.1250}}
  >{\raggedright\arraybackslash}p{(\linewidth - 14\tabcolsep) * \real{0.1250}}
  >{\raggedright\arraybackslash}p{(\linewidth - 14\tabcolsep) * \real{0.1250}}
  >{\raggedright\arraybackslash}p{(\linewidth - 14\tabcolsep) * \real{0.1250}}
  >{\raggedright\arraybackslash}p{(\linewidth - 14\tabcolsep) * \real{0.1250}}
  >{\raggedright\arraybackslash}p{(\linewidth - 14\tabcolsep) * \real{0.1250}}
  >{\raggedright\arraybackslash}p{(\linewidth - 14\tabcolsep) * \real{0.1250}}@{}}
\toprule\noalign{}
\begin{minipage}[b]{\linewidth}\raggedright
Place
\end{minipage} & \begin{minipage}[b]{\linewidth}\raggedright
Rows
\end{minipage} & \begin{minipage}[b]{\linewidth}\raggedright
Pieces
\end{minipage} & \begin{minipage}[b]{\linewidth}\raggedright
Models, code / generated
\end{minipage} & \begin{minipage}[b]{\linewidth}\raggedright
Cloud
\end{minipage} & \begin{minipage}[b]{\linewidth}\raggedright
Pictures
\end{minipage} & \begin{minipage}[b]{\linewidth}\raggedright
Faults before / after
\end{minipage} & \begin{minipage}[b]{\linewidth}\raggedright
Resting
\end{minipage} \\
\midrule\noalign{}
\endhead
\bottomrule\noalign{}
\endlastfoot
Hub & 74 & 460 & 88 / 12 & €3.65 & \$1.07 & not counted & 405 of 460 \\
Lab & 47 & 412 & 54 / 22 & €8.36 & \$4.15 & 13 / 12 & 373 of 412 \\
Greenhouse & 59 & 1,248 & 146 / 3 & €3.07 & \$2.13 & 7 / 3 & 1,154 of
1,248 \\
Workshop & 54 & 814 & 97 / 8 & €2.18 & \$3.35 & 3 / 9 & 671 of 814 \\
Airlock & 33 & 330 & 73 / 1 & €2.53 & \$0.13 & 1 / 5 & 296 of 330 \\
Walkway tube & 13 & 39 & 23 / 0 & €0.81 & \$0 & 0 / 2 & 38 of 39 \\
Habitat & 53 & 637 & 100 / 2 & €5.70 & \$1.34 & 1 / 4 & 543 of 637 \\
Garage & 41 & 602 & 75 / 19 & €4.29 & \$8.71 with the hangar & 3 / not
counted, with the hangar & 493 of 602 \\
Hangar & 37 & 1,686 & 61 / 13 & €4.22 & with the garage & with the
garage & 865 of 1,686 \\
Wreck & 14 & 22 & 3 / 13 & €4.88 & \$2.68 & 2 / 1 & 22 of 22 \\
Old station & 17 & 28 & 10 / 14 & €4.93 & \$3.22 & 5 / 2 & 29 of 29 \\
Mars camp & 33 & 218 & 51 / 5 & €3.47 & \$1.21 & 4 / 8 & 127 of 218 \\
Launch view & 4 & 67 & 4 / 1 & €0.33 & \$0.40 & 1 / 2 & 8 of 67 \\
Square & 40 & 4,135 & 74 / 2 & €1.64 & \$1.07 & 3 / 4 & 3,712 of
4,135 \\
Street & 33 & 1,198 & 45 / 2 & €0.98 & \$1.74 & 7 / 6 & 1,179 of
1,198 \\
Stairwell & 21 & 118 & 42 / 1 & €1.23 & \$2.41 with the flat & 2 / 6 &
117 of 118 \\
Flat & 38 & 52 & 35 / 9 & €2.44 & with the stairwell & 5 / 6 & 50 of
52 \\
World 1 & 611 & 12,066 & 981 / 127 & €54.71 & & & 10,082 of 12,066 \\
\end{longtable}
}

\section{Development history}\label{development-history}

The framework grew out of building the game 2099 between late September
and early October 2026, with one owner directing coding agents. In its
first form it joined four parts: rooms laid out in code, a scene plan
from a world model that painted depth panoramas of those rooms, props
from a picture-to-3D pipeline, and coding agents that fitted each room
to its plan. Eight places were fitted out to their plans in a single
day. The owner found that their layouts were close to the plans, but
that their styles were mixed, with models from earlier rounds and with
objects that the plans never showed. The cause lay in the workflow:
every pass added things to a room, and no pass checked what had been
placed against the plan. This led to the inventory as a strict list,
with nothing made or placed outside it, and to deciding the look before
any money is spent. The structure of each room then came from concept
cutaways drawn from the owner's own reference pictures; for the hub he
picked one cutaway out of twenty. Scale moved into a dimensioned plan
drawn before the concept, after the first cutaway came out at
doll's-house scale: about 13 m across, judged by the people drawn in it,
against the 8 to 9 m that had been asked for.

\subsection{The hub in six rounds}\label{the-hub-in-six-rounds}

The mechanisms of Section 2 then came out of six rounds of work on one
room, the central hub of the game 2099, on 6 and 7 October 2026. Each
round answered the creator's verdict on the round before (Table C1). In
the first two rounds, the creator chose between the old and the new
method by comparing them side by side from the same cameras.

\textbf{Table C1.} The hub's six rounds. Cost is given as rented GPU
time in euros and picture-model calls in dollars. Frame time is the mean
frame time on an RTX 5080, in the night scene, at 1600×900.

{\def\LTcaptype{none} % do not increment counter
\begin{longtable}[]{@{}
  >{\raggedright\arraybackslash}p{(\linewidth - 8\tabcolsep) * \real{0.2000}}
  >{\raggedright\arraybackslash}p{(\linewidth - 8\tabcolsep) * \real{0.2000}}
  >{\raggedright\arraybackslash}p{(\linewidth - 8\tabcolsep) * \real{0.2000}}
  >{\raggedright\arraybackslash}p{(\linewidth - 8\tabcolsep) * \real{0.2000}}
  >{\raggedright\arraybackslash}p{(\linewidth - 8\tabcolsep) * \real{0.2000}}@{}}
\toprule\noalign{}
\begin{minipage}[b]{\linewidth}\raggedright
Round
\end{minipage} & \begin{minipage}[b]{\linewidth}\raggedright
Change
\end{minipage} & \begin{minipage}[b]{\linewidth}\raggedright
Creator's verdict
\end{minipage} & \begin{minipage}[b]{\linewidth}\raggedright
Cost
\end{minipage} & \begin{minipage}[b]{\linewidth}\raggedright
Frame time
\end{minipage} \\
\midrule\noalign{}
\endhead
\bottomrule\noalign{}
\endlastfoot
1, two walls, A/B & surface library, routing by shape, room checks &
chose the new method (``Y looks much cleaner'') & €0.88 & 7.81 ms vs
7.53 \\
2, two walls, A/B & 71-variant library, shared texture sets, roof check
& chose the new method again & €1.70 & 8.05 ms vs 8.24 \\
3, whole room & all stages end to end, 74 models & old models were still
loaded; labels and floor fittings were wrong & €1.67 & 6.78 ms (old
9.22) \\
4, whole room & detail through the pipeline, picture colours kept &
``pretty much all of this became messy'' & €8.50, \$5.09 & 9.6 ms \\
5, eight pieces, A/B & shape only from the pipeline, library surfaces,
parts check & new method preferred for the door, notice board and
porthole; three pieces still wrong & €0.21 & not run \\
6, whole room & round 5 everywhere, one model per object, straightness
check & ``this looks amazing, I have no negative feedback'' & €3.48,
\$1.07 & 5.6 ms \\
\end{longtable}
}

Three findings shaped the method. First, deterministic checks catch what
agents miss. In round one, the share of rays that leaked out through
walls and roof fell from 1.53 percent of 5,817 rays to zero, and the
colour noise within surfaces fell from a median of 31.1 to 3.2 percent.
Only walking through the loaded room, however, showed that 26 old models
were still in it, and that finding became the made-only check. Second,
when each generated piece kept the colours of its own picture, every
piece brought its own wear, and the room looked inconsistent; surfaces
have since come only from the library. Third, generated composite
objects melted their parts together, and generated boxes leaned (one
locker by 13 degrees). This led to one model per real-world object and
to the straightness check (Figure 4b). By round six the room used 0.45
million triangles, against 1.96 million in round four, and 187 MB on
disk, against 487 MB.

\section{Check details}\label{check-details}

\begin{itemize}
\tightlist
\item
  \textbf{Room checks} take about two seconds per room and run on the
  laid-out meshes. They cast rays to find light leaking out through
  walls and roof, find pieces placed in front of a real opening, find
  roof and wall fittings more than 3 cm off their surface or turned more
  than 5 degrees from it, and count pieces, triangles and lamps against
  the room's budget.
\item
  \textbf{Door check:} every door must stand in a wall or partition that
  fully separates its two sides. The check tests this by walking a grid
  from one side to the other.
\item
  \textbf{Concept-density check:} the concept picture is cut into crops,
  and every visible element must have an inventory row or a written
  reason for leaving it out. After the build, the place is rendered from
  the concept's camera and shown beside the concept.
\item
  \textbf{Parts check:} in the code build, every part shown in the
  close-up picture must be at least 10 percent visible from the front,
  and no label may cover a part.
\item
  \textbf{Model check:} every model must be a closed solid with walls at
  least 3 mm thick. A generated box fails if one of its sides tilts more
  than 2 degrees or is warped. The step that makes a generated model
  solid stops itself if the model's bounding box grows more than 3
  percent.
\item
  \textbf{Palette sweep and storage budget:} the colours of every model
  are compared with the place's library palette, and each room may use
  at most 250 MB on disk and 640 MB of textures on the graphics card.
\item
  \textbf{Scene check and made-only check:} in the engine, these checks
  find anything that floats, sinks or overlaps, and any visible mesh
  that the framework did not make. A made-only finding can never be
  waived.
\item
  \textbf{Close-up shape check:} the picture must have a plain border
  all round, one object filling a sensible share of the frame, and an
  outline within a set proportion of the inventory row's box; the
  judge's seven questions must pass by a majority of three sampled
  answers. A close-up passes only when both halves pass.
\item
  \textbf{Settling and resting check:} on the OpenUSD stage, loose
  objects are dropped with their own triangles onto the ground and the
  fixed objects. An object that would turn more than 15 degrees or drift
  more than 30 cm keeps its laid pose and is marked as not resting;
  small debris, with no side over 1.4 m, may turn freely and drift up to
  1 m. The resting check then fails anything floating, tipping, sunk
  below its contact points, lying inside another piece or marked as not
  resting. Hung objects are exempt by their inventory anchor, and fixed
  objects are not judged for tipping.
\item
  \textbf{Resting triage:} each object's support comes from its
  inventory row: a glowing part or screen is judged with its host; a
  fixed row stands as laid but may not float unless joined to another
  piece; a kit's shell pieces stand in the kit's frame; a hanging row
  and a wall or ceiling row need one contact with anything; a floor row
  and a row standing on another need four points of their footing (the
  points within 10 cm of their lowest) within 3 cm of a surface under
  them, with their weight within 5 cm of the area those points span. An
  overlap deeper than 3 cm counts only if the two still meet after the
  loose one moves 5 mm away. Cables may sink 6 cm, cloths 8 cm and
  ground debris 30 cm. From a drop run, an object that moves more than
  0.2 m or turns more than 8 degrees is a real fault.
\item
  \textbf{Completion gate:} the requirements given in Section 2 under
  the canonical scene, each with its result tied to a fingerprint of the
  stage's files and inputs; a result made on another stage counts as
  unknown.
\item
  \textbf{Pick lock and render-and-compare:} a picked row's objects must
  show the model with the picked hash, or the row must carry the
  creator's written replacement. Each object with a close-up is rendered
  from eight directions 20 degrees up, and the direction with the
  highest DINOv2 likeness stands for the close-up's camera; the outline
  score is the overlap of the two silhouettes, each cut to its box.
\item
  \textbf{Brightness check:} each view of a place's stage on its review
  page is shown beside the game's shot from the same spot, and a view
  whose mean brightness is under half the game's fails.
\end{itemize}

\section{Models and licences}\label{models-and-licences}

Every model in the framework must allow commercial use of its output
(Table E1). Candidate models were dropped under this rule; among them
were NVIDIA's Lyra 2.0, whose weights are licensed for internal research
only, and Hunyuan3D 2.1, whose licence does not apply in the European
Union. We read each licence for its clauses on commercial use. Four of
the models read pictures through Meta's DINO family of models, each
through its own copy. What passes from one model to the next is always
pictures, masks, depth maps or meshes, never one model's internal
features. The characters' licence check found that GarmentCode's fork of
NVIDIA Warp, which drapes its patterns, allows non-commercial research
only, so Blender's cloth simulation replaced it. The one open doubt
among the characters' models is the face LoRA, whose training data is
not disclosed.

\textbf{Table E1.} Models in the framework.

{\def\LTcaptype{none} % do not increment counter
\begin{longtable}[]{@{}
  >{\raggedright\arraybackslash}p{(\linewidth - 4\tabcolsep) * \real{0.3333}}
  >{\raggedright\arraybackslash}p{(\linewidth - 4\tabcolsep) * \real{0.3333}}
  >{\raggedright\arraybackslash}p{(\linewidth - 4\tabcolsep) * \real{0.3333}}@{}}
\toprule\noalign{}
\begin{minipage}[b]{\linewidth}\raggedright
Model
\end{minipage} & \begin{minipage}[b]{\linewidth}\raggedright
Role
\end{minipage} & \begin{minipage}[b]{\linewidth}\raggedright
Licence
\end{minipage} \\
\midrule\noalign{}
\endhead
\bottomrule\noalign{}
\endlastfoot
Nano Banana Pro {[}30{]} & concepts, close-ups the open model fails,
ground skins & paid API \\
Qwen-Image-Edit-2511 {[}28{]} & close-ups, first & Apache-2.0 \\
Qwen3.8-27B {[}29{]} & the close-ups' shape check; each part's material
& Apache-2.0 \\
Pixal3D {[}17{]} & picture to 3D & MIT; DINOv3 licence {[}44{]} for its
encoder \\
PartCrafter {[}21{]} & parts & MIT; its non-commercial background
remover not used \\
SegviGen {[}45{]} & parts of our own model (trial) & MIT, on TRELLIS.2
(MIT); DINOv3 licence for its encoder; its non-commercial background
remover and nvdiffrast not used \\
GeoSAM2 {[}46{]} & parts of our own model, seeded by the close-up's
regions (trial) & Apache-2.0 \\
Depth Anything V2 Small {[}20{]} & ground relief & Apache-2.0 \\
ProcFunc, Infinigen shaders & surfaces & BSD-3-Clause \\
MOSS-SoundEffect v2.0, CLAP & sound & Apache-2.0 \\
World Labs Marble {[}47{]} & optional world step & paid; outputs owned
by paid users \\
FLUX.2 klein 4B {[}35{]} & earlier prop pictures; the characters' A-pose
pictures from words and their hair's clay pictures & Apache-2.0 \\
MoGe-2 {[}19{]}, SAM 3 {[}18{]} & measuring and finding objects in
pictures & MIT; SAM Licence \\
Kimodo {[}23{]} & the characters' motion from sentences & Apache-2.0
code; NVIDIA Open Model License for the Kimodo-SOMA-RP-v1.1 weights,
commercial use allowed \\
LLM2Vec {[}48{]} on Meta-Llama-3-8B-Instruct {[}49{]} & Kimodo's text
encoder, used only while clips are made; nothing from it ships & Llama 3
Community License: commercial use allowed, the 700 million monthly users
clause does not apply, ``Built with Meta Llama 3'' shown, no exclusion
in the European Union for the text-only 8B model \\
SAM 3D Body {[}22{]}, SAM3DBody-cpp {[}36{]} & body shape from a picture
& SAM Licence for the weights, SAM3DBody-cpp's converted ONNX and GGUF
weights included whatever their tag says; MIT for SAM3DBody-cpp's
code \\
SOMA-X and MHR & the characters' bodies & Apache-2.0 \\
GarmentCode {[}37{]}, Blender's cloth {[}38{]} & the clothes' sewing
patterns and their drape & MIT; Blender is GPL and its output is ours;
GarmentCode's Warp fork (non-commercial research only) not used \\
Hi3DGen {[}39{]} & the characters' hair & MIT code; MIT and Apache-2.0
weights \\
FLUX.2 klein base 4B {[}40{]} with the refcontrol reference-depth LoRA
{[}41{]} & the characters' faces & Apache-2.0, both; the LoRA's training
data is not disclosed \\
UniRig {[}42{]} & the skeleton and skin of a quadruped & MIT code and
weights \\
MakeHuman eyes & the characters' eyes & CC0 \\
\end{longtable}
}

On our tier, Nano Banana Pro allows 250 pictures a day and 20 a minute;
the cheaper Nano Banana (gemini-2.5-flash-image) allows 2,000 a day and
500 a minute. We tested whether the cheaper model could draw the
close-ups instead, on the same 10 lab objects with the same prompts and
inputs (Table E2). It refused the standard prompt for 6 of the 10
objects and drew those only from the crop alone, and its pictures were
about one megapixel with the object small in the frame. Close-ups
therefore stay on Nano Banana Pro.

\textbf{Table E2.} Ten lab close-ups drawn by Nano Banana instead of
Nano Banana Pro.

{\def\LTcaptype{none} % do not increment counter
\begin{longtable}[]{@{}ll@{}}
\toprule\noalign{}
Result & Objects \\
\midrule\noalign{}
\endhead
\bottomrule\noalign{}
\endlastfoot
good & 2 \\
usable & 3 \\
marginal & 1 \\
wrong shape (microscope, chair, sample tray) & 3 \\
refused (glovebox) & 1 \\
\end{longtable}
}

The close-ups' shape check was tuned on 40 close-ups that four open
picture models drew of the lab's ten objects, each scored by hand as
good, usable, marginal, wrong or failed, and on Nano Banana Pro's ten
pictures of the same objects. With the majority of three sampled
answers, the check agreed with the hand score on 29 of the 40, passed
all ten of the closed model's pictures and two marginal ones, passed
none of the pictures scored wrong or failed, and failed nine that were
good or usable, each of which then cost a picture from the closed model.
The judge's single answer at temperature zero was rejected: asked the
same 51 questions twice in one batch, it gave a different verdict on 7
of them. On the lab's 22 generated rows, an earlier version of the check
accepted 8 of the open model's close-ups and sent 14 to the closed
model; 7 of the closed model's first pictures failed for drawing the
room or dimension lines around the object, and all 7 passed when redrawn
from the crop alone. On the 17 rows that world 1 still needed
afterwards, from three places, none of the open model's pictures passed;
the closed model's pictures were accepted for 10 rows, and 7 rows were
left for the creator to look at.

\textbf{Table E3.} The close-up stage on world 1. Cloud is rented card
time in euros; the closed model's pictures are in dollars.

{\def\LTcaptype{none} % do not increment counter
\begin{longtable}[]{@{}
  >{\raggedright\arraybackslash}p{(\linewidth - 12\tabcolsep) * \real{0.1429}}
  >{\raggedright\arraybackslash}p{(\linewidth - 12\tabcolsep) * \real{0.1429}}
  >{\raggedright\arraybackslash}p{(\linewidth - 12\tabcolsep) * \real{0.1429}}
  >{\raggedright\arraybackslash}p{(\linewidth - 12\tabcolsep) * \real{0.1429}}
  >{\raggedright\arraybackslash}p{(\linewidth - 12\tabcolsep) * \real{0.1429}}
  >{\raggedright\arraybackslash}p{(\linewidth - 12\tabcolsep) * \real{0.1429}}
  >{\raggedright\arraybackslash}p{(\linewidth - 12\tabcolsep) * \real{0.1429}}@{}}
\toprule\noalign{}
\begin{minipage}[b]{\linewidth}\raggedright
Run
\end{minipage} & \begin{minipage}[b]{\linewidth}\raggedright
Close-ups
\end{minipage} & \begin{minipage}[b]{\linewidth}\raggedright
Open model accepted
\end{minipage} & \begin{minipage}[b]{\linewidth}\raggedright
Closed model accepted
\end{minipage} & \begin{minipage}[b]{\linewidth}\raggedright
Left for the creator
\end{minipage} & \begin{minipage}[b]{\linewidth}\raggedright
Cloud
\end{minipage} & \begin{minipage}[b]{\linewidth}\raggedright
Closed model
\end{minipage} \\
\midrule\noalign{}
\endhead
\bottomrule\noalign{}
\endlastfoot
Lab, 22 rows (earlier check: one answer) & 22 & 8 & 14 & 0 & €4.12 &
\$2.81 \\
World 1's remaining rows (three places) & 17 & 0 & 10 & 7 & €3.06 &
\$4.02 \\
\end{longtable}
}

The picture-to-3D step, the part splitting, the surface bakes, the sound
generation, all checks and the assembly run automatically. Coding agents
(Claude Code {[}50{]}) write the plans, the inventories and every code
builder; in the hub's final round, 86 of its 98 models were built in
code. Some steps were still done by hand: some generated objects were
turned to face the right way by eye, and a few labels were placed by
hand before a rule took over that job. The concepts and the verdicts are
the creator's own choices.

\section{Surfaces, sound and light}\label{surfaces-sound-and-light}

The surface library holds families of variants: painted, steel,
aluminium, rubber, plastic, cable, fabric, composite, glass, screen,
light and print. After the hub's second round it held 71 variants made
from 14 ProcFunc recipes. Each variant is a recipe with settings, takes
its colours only from the palette, and is baked into base-colour,
metal-roughness and normal maps at a texture density that depends on how
close the player can get. Wear has three levels and comes from causes:
edges, and kicks within 32 cm of the floor that a piece stands on. Every
sound starts from a short written brief about its object. Four takes are
generated for each sound, scored by how well CLAP matches them to their
prompt minus any measured faults (clipping, unsteady loops or bad
joins), and levelled to one loudness per kind of sound under a true peak
of -3 dBTP. A creator may swap any take. For the game 2099, all 61
sounds were generated as 244 takes in one batch for €1.61, and the
creator chose to keep earlier recordings for three kinds of sound:
breathing, a sliding door, and footsteps on steel.

Light adds up, so it is planned as layers. For each kind of module, the
bounce light of each group of lamps, and of a few sun directions through
the windows, will be baked on rented machines as a separate layer, and
the engine will mix these layers live according to each lamp's
brightness. Outdoor light and the light on moving objects stay live.
This part is not built yet.

\section{Cost and infrastructure}\label{cost-and-infrastructure}

Every heavy step runs on rented machines that are deleted after each
batch: NVIDIA graphics cards for picture-to-3D, part splitting and
baking, and processor-only machines when no cards are in stock. The
engine's tests, its screenshots and its frame-time benchmark also run on
rented machines. A full test run took 542 s there, against 596 s on the
local computer, and frame times measured on an L4 card are converted to
the local RTX 5080 by a measured factor of 3.47. A picture costs about
\$0.134 at the resolution we use. The picture model's limits on speed
are described under Limitations.

A batch asks for a class of machine, such as one card with 24 GB, and
takes whatever the provider has in stock across three zones and three
card types, because on 8 October the cheapest type was out of stock in
one zone for most of the day. Table G1 compares the same jobs on two
card types. An H100 runs ten Pixal3D takes at once in 31 GB of its 80
GB, and makes nearly three times as many models an hour as an L4. It has
no ray-tracing cores, so it bakes 2.6 times slower than an L4. Pixal3D
therefore takes 80 GB cards first, and a bake takes an H100 only when no
other card has been free for five minutes. The L40S was out of stock on
both attempts, and the provider's P100 machines started but their driver
did not see the card.

\textbf{Table G1.} The same jobs on two card types, 8 October 2026. Wait
is the time from the order until the machine answers; machine time is
the whole rental, setup included.

{\def\LTcaptype{none} % do not increment counter
\begin{longtable}[]{@{}
  >{\raggedright\arraybackslash}p{(\linewidth - 12\tabcolsep) * \real{0.1429}}
  >{\raggedright\arraybackslash}p{(\linewidth - 12\tabcolsep) * \real{0.1429}}
  >{\raggedright\arraybackslash}p{(\linewidth - 12\tabcolsep) * \real{0.1429}}
  >{\raggedright\arraybackslash}p{(\linewidth - 12\tabcolsep) * \real{0.1429}}
  >{\raggedright\arraybackslash}p{(\linewidth - 12\tabcolsep) * \real{0.1429}}
  >{\raggedright\arraybackslash}p{(\linewidth - 12\tabcolsep) * \real{0.1429}}
  >{\raggedright\arraybackslash}p{(\linewidth - 12\tabcolsep) * \real{0.1429}}@{}}
\toprule\noalign{}
\begin{minipage}[b]{\linewidth}\raggedright
Card
\end{minipage} & \begin{minipage}[b]{\linewidth}\raggedright
Job
\end{minipage} & \begin{minipage}[b]{\linewidth}\raggedright
At once
\end{minipage} & \begin{minipage}[b]{\linewidth}\raggedright
Wait
\end{minipage} & \begin{minipage}[b]{\linewidth}\raggedright
Machine time
\end{minipage} & \begin{minipage}[b]{\linewidth}\raggedright
Cost
\end{minipage} & \begin{minipage}[b]{\linewidth}\raggedright
Per take
\end{minipage} \\
\midrule\noalign{}
\endhead
\bottomrule\noalign{}
\endlastfoot
L4, 24 GB & Pixal3D, 3 close-ups & 3 & 0.4 min & 28.0 min & €0.38 &
€0.13, 6.4 an hour \\
H100, 80 GB & Pixal3D, the same 3 & 3 & 0.6 min & 13.9 min & €0.67 &
€0.22, 12.9 an hour \\
H100, 80 GB & Pixal3D, 10 close-ups & 10 & 0.6 min & 33.1 min & €1.62 &
€0.16, 18.1 an hour \\
L4, 24 GB & bake: 60 surfaces at 3 wear levels & 1 & 0.4 min & 3.8 min
(bake 138 s) & €0.05 & \\
H100, 80 GB & the same bake & 1 & 0.6 min & 7.5 min (bake 360 s) & €0.38
& \\
\end{longtable}
}

Since 8 October, every kind of batch is spread over many machines at
once: each machine is set up once and takes the next job from one queue
as it finishes one, and a batch rents enough machines that each works
about as long as its setup takes. Table G2 compares the same batches on
one machine and spread. A spread batch ended 1.2 to 2.7 times sooner and
cost 4 to 63 percent more, because every added machine spends part of
its time setting up.

\textbf{Table G2.} The same batches on one machine and spread over
several, 8 October 2026.

{\def\LTcaptype{none} % do not increment counter
\begin{longtable}[]{@{}
  >{\raggedright\arraybackslash}p{(\linewidth - 8\tabcolsep) * \real{0.2000}}
  >{\raggedright\arraybackslash}p{(\linewidth - 8\tabcolsep) * \real{0.2000}}
  >{\raggedright\arraybackslash}p{(\linewidth - 8\tabcolsep) * \real{0.2000}}
  >{\raggedright\arraybackslash}p{(\linewidth - 8\tabcolsep) * \real{0.2000}}
  >{\raggedright\arraybackslash}p{(\linewidth - 8\tabcolsep) * \real{0.2000}}@{}}
\toprule\noalign{}
\begin{minipage}[b]{\linewidth}\raggedright
Batch
\end{minipage} & \begin{minipage}[b]{\linewidth}\raggedright
Machines
\end{minipage} & \begin{minipage}[b]{\linewidth}\raggedright
Time
\end{minipage} & \begin{minipage}[b]{\linewidth}\raggedright
Cost
\end{minipage} & \begin{minipage}[b]{\linewidth}\raggedright
On one machine
\end{minipage} \\
\midrule\noalign{}
\endhead
\bottomrule\noalign{}
\endlastfoot
Blender, 6 jobs & 3 L4 & 7.0 min & €0.28 & 14.3 min, €0.20 \\
Library bake, 22 jobs & 4 L4 & 10 min & €0.43 & 27 min, €0.35 \\
Pictures (FLUX.2 klein), 120 & 3 L4 & 5 min & €0.20 & 10 min, €0.14 \\
Pictures without light, 32 & 2 L4 & 4 min & €0.10 & 6 min, €0.08 \\
PartCrafter, 12 runs & 2 L4 & 19.3 min & €0.49 & 35.2 min, €0.47 \\
Shape-check judge, 184 questions & 2 H100 & 16 min & €1.48 & 19 min,
€0.91 \\
Qwen-Image-Edit, 16 close-ups & 3 H100 & 11 min & €1.53 & 26 min,
€1.24 \\
\end{longtable}
}

Table G3 sums the cloud ledger, in which every runner records each
machine it rents, by type of machine, from the first batch on 29
September to the end of 8 October 2026. It holds the game 2099's own
test and benchmark machines as well as the framework's batches, and the
machines of the earliest batches, worth €34.69, were recorded without
their type.

\textbf{Table G3.} Rented machines by type, 29 September to 8 October
2026.

{\def\LTcaptype{none} % do not increment counter
\begin{longtable}[]{@{}llll@{}}
\toprule\noalign{}
Machine & Rentals & Machine time & Cost \\
\midrule\noalign{}
\endhead
\bottomrule\noalign{}
\endlastfoot
L4, 24 GB & 793 & 331.1 h & €266.19 \\
H100 PCIe, 80 GB & 29 & 7.4 h & €21.69 \\
two H100 SXM, 80 GB each & 14 & 1.5 h & €11.03 \\
32-core processor, 128 GB & 26 & 5.4 h & €9.85 \\
two L4, 24 GB each & 3 & 3.4 h & €5.38 \\
16-core processor, 64 GB & 13 & 2.3 h & €4.81 \\
32-core processor, 64 GB & 5 & 0.9 h & €4.26 \\
L40S, 48 GB & 6 & 2.7 h & €4.07 \\
P100, 16 GB & 2 & 0.1 h & €2.44 \\
16-core processor, 128 GB & 2 & 1.6 h & €1.36 \\
type not recorded & & & €34.69 \\
total & & & €366.76 \\
\end{longtable}
}

The character maker runs a person's whole chain on one card (Table G4).
Its machine is set up once with six environments side by side, after the
CUDA 12.6 toolkit that SAM3DBody-cpp needs: 11.0 minutes on an H100 and
21 to 25 minutes on an L4. On an H100, SAM3DBody-cpp's network takes 1.3
s and its whole call about 4 s once warm, 44 s cold. Against the Python
SAM 3D Body on the same picture, its shape parameters correlate at 0.98,
and its body at rest is 23 mm shorter (1.548 m against 1.571 m), with a
mean gap of 17.6 mm between vertices and joints 19 mm apart on average
and 25 mm at worst. The difference comes from the person box:
SAM3DBody-cpp finds its own with YOLO, while the Python run was given
the whole picture. Kimodo with its text encoder on the same card writes
a clip in 30 s once warm, and the first in 82 s, which includes
downloading the Llama 3 encoder. Nev's 19 clips had been made already,
so her rebuild wrote none. On her walk, the built character's skinning,
with 27 joints and thinned weights as an engine plays it, keeps within
6.8 mm of the body model's own posed points on average and 52.6 mm at
worst (4.9 mm on average off the hands), and its file is 5.05 MB. Her
far body has 612 vertices and 1,220 triangles, her work suit 31,260
triangles and her space suit 44,846. The development machine on which
the chain was built and Nev was rebuilt ran for 163 billed minutes and
cost €7.79 in all. On an L4, the same chain made the player from the
creator's drawing of him, on his own body 1.742 m tall, in 26.2 minutes,
for about €0.35 of the card at €0.79 an hour. The face is the slow step
on an L4: klein base 4B with the LoRA, 50 steps and two seeds, took 13.6
of those minutes against 2.7 on the H100. The L4 was set up in 24.8
minutes. A new person made from words alone, an engineer, first reached
the build on a second L4 and failed at the arm stripe of her space suit,
because the drape had left her upper arm bare; the drape now tries again
when a sleeve does not cover the upper arm. The batch of both people
cost €1.50 for 60 minutes on three L4 machines, one of which failed its
setup after 25 minutes because Kimodo's C++ extension was built while
the CUDA toolkit installed beside it; the toolkit now installs first.
Made again on an H100, the engineer took 20.9 minutes: klein drew her
picture in 0.27 minutes, the drapes took 6.65 and the review renders
5.96, because a fresh H100 first compiles Cycles' kernels. On a 16-core
processor machine, Blender drapes the work suit in 1.5 minutes and the
space suit in 0.6 minutes, with every seam within 3 cm and no cloth
inside the body.

The animals' route rents cheaper cards. One Pixal3D run on an L40S made
both animals' meshes in 11.2 minutes for €0.29, and UniRig rigs a model
in about 70 s on an L4. The rig, the clips, the far levels and the
review renders take 5.2 minutes of an L4 for €0.08. A processor machine
costs 10 to 15 times more for the same job: the fish's own Blender run
went to a 32-core processor machine when no card was free and cost
€0.85. UsdSkel skins the animals to within 0.2 µm of their glTF skinning
for the dog and 0.01 µm for the fish. The fish's levels have 40,000,
6,000 and 532 triangles, the dog's 39,999, 5,999 and 1,500.

\textbf{Table G4.} Characters made by the character maker, 9 October
2026. Minutes of card time per group of steps; a person's cost is the
card time of the chain alone, without the setup; the fish and the dog
shared one Pixal3D run on an L40S, split evenly here, and their costs
are on cards for every step.

{\def\LTcaptype{none} % do not increment counter
\begin{longtable}[]{@{}
  >{\raggedright\arraybackslash}p{(\linewidth - 20\tabcolsep) * \real{0.0909}}
  >{\raggedright\arraybackslash}p{(\linewidth - 20\tabcolsep) * \real{0.0909}}
  >{\raggedright\arraybackslash}p{(\linewidth - 20\tabcolsep) * \real{0.0909}}
  >{\raggedright\arraybackslash}p{(\linewidth - 20\tabcolsep) * \real{0.0909}}
  >{\raggedright\arraybackslash}p{(\linewidth - 20\tabcolsep) * \real{0.0909}}
  >{\raggedright\arraybackslash}p{(\linewidth - 20\tabcolsep) * \real{0.0909}}
  >{\raggedright\arraybackslash}p{(\linewidth - 20\tabcolsep) * \real{0.0909}}
  >{\raggedright\arraybackslash}p{(\linewidth - 20\tabcolsep) * \real{0.0909}}
  >{\raggedright\arraybackslash}p{(\linewidth - 20\tabcolsep) * \real{0.0909}}
  >{\raggedright\arraybackslash}p{(\linewidth - 20\tabcolsep) * \real{0.0909}}
  >{\raggedright\arraybackslash}p{(\linewidth - 20\tabcolsep) * \real{0.0909}}@{}}
\toprule\noalign{}
\begin{minipage}[b]{\linewidth}\raggedright
Character
\end{minipage} & \begin{minipage}[b]{\linewidth}\raggedright
Card
\end{minipage} & \begin{minipage}[b]{\linewidth}\raggedright
Picture
\end{minipage} & \begin{minipage}[b]{\linewidth}\raggedright
Body or mesh
\end{minipage} & \begin{minipage}[b]{\linewidth}\raggedright
Rig and clips
\end{minipage} & \begin{minipage}[b]{\linewidth}\raggedright
Clothes
\end{minipage} & \begin{minipage}[b]{\linewidth}\raggedright
Head, hair and face
\end{minipage} & \begin{minipage}[b]{\linewidth}\raggedright
Build and review
\end{minipage} & \begin{minipage}[b]{\linewidth}\raggedright
Total
\end{minipage} & \begin{minipage}[b]{\linewidth}\raggedright
Cost
\end{minipage} & \begin{minipage}[b]{\linewidth}\raggedright
Setup
\end{minipage} \\
\midrule\noalign{}
\endhead
\bottomrule\noalign{}
\endlastfoot
Nev, from her picture & H100 PCIe & 0 (her own) & 0.13 & 0 (made
already) & 4.75 & 4.23 & 2.09 & 11.2 min & €0.54 & not measured
cleanly \\
The engineer, from words & H100 PCIe & 0.27 & 0.91 & 0 (made already) &
6.65 & 5.21 & 7.84 & 20.9 min & €1.00 & 11.0 min \\
The player, from the creator's drawing & L4 & 0 (the creator's drawing)
& 0.34 & 0 (made already) & 5.18 & 17.08 & 3.60 & 26.2 min & €0.35 &
24.8 min \\
Dog, from words & L4; L40S for the mesh & 2 (4 pictures) & 5.6 & about
1.2 (UniRig) & & & 5.2 & about 14 min & about €0.40 & about 1.5 min
(UniRig) \\
Fish, from a close-up & L40S for the mesh; L4 & 0 (given) & 5.6 & in
code & & & on an L4 as the dog's & & about €0.32 & \\
\end{longtable}
}

\begin{center}\fbox{\begin{minipage}{0.94\linewidth}\setlength{\parskip}{4pt}\small

\textbf{Gap: cost tables.} These tables will give the cost per kind of
batch and per world, reconciled with the cloud provider's bill. They
wait on the worlds and on the bill; the figures per place so far
(Appendix B) are lower bounds taken from the build logs.

\end{minipage}}\end{center}

\end{document}
